\chapter{$\mathrm{GL}_n$ 的 \texorpdfstring{$L$}{L} 函数解析理论}\label{ch:09}

\begin{center}
J.\,W. Cogdell
\end{center}

本章的目的是描述整体域上 $\GL_n$ 的尖点自守表示的 $L$ 函数的解析理论。
$\GL_n$ 的 $L$ 函数有两条进路：一条经由积分表示，另一条经由 Eisenstein 级数的
Fourier 系数分析。本章讨论的是积分表示这一条。

自守形式（或模形式）$L$ 函数的积分表示理论，起源于 Riemann 关于 $\zeta$ 函数的论文。
然而这套理论真正发展成形，是在 $\mathrm{SL}_2(\mathbb{Z})$ 同余子群的模形式
$L$ 函数这一经典框架内，由 Hecke 及其学派完成的。
我们今日理论的很大一部分，直接由 Hecke 的工作生长而来。
自守表示的 $L$ 函数最早由 Jacquet 与 Langlands 对 $\GL_2$ 建立，
其进路遵循 Hecke，并融入 Tate 论文的局部—整体技巧。
$\GL_n$ 的理论随后沿同一路线发展，散见于 Jacquet、Piatetski-Shapiro 与 Shalika
以各种组合写成的长系列论文。
Hecke 除给自守形式配上 $L$ 函数外，还给出一个判别法，判定 Dirichlet 级数是否来自模形式，
即所谓 Hecke 逆定理。
在自守表示的框架内，$\GL_2$ 的逆定理由 Jacquet 与 Langlands 发展，
由 Jacquet、Piatetski-Shapiro 与 Shalika 推广并显著加强到 $\GL_3$，
随后又与 Piatetski-Shapiro 一起推到 $\GL_n$。
我们在此试图呈现的，是这一工作的提要；本章的扩充版本见 [1]。
\begin{annote}
Hecke 的逆定理是本书若干章的共同源头，值得先说清它的逻辑形状。
正向的对应把模形式送到 Dirichlet 级数，靠的是 Mellin 变换；
逆向则把这个积分「反解」出来，用函数方程与整性把 $f$ 重新造回。
这条思路在 $\GL_n$ 上仍然奏效，只是 Mellin 变换要换成第 2 节的 Euler 积分，
而「反解」要换成对 Fourier 展开逐项作 Mellin 逆变换。
Cogdell 在 [1] 的长篇讲义中把这条主线写得很完整，本章只取提要。
另需留意，本章的 $L$ 函数一律取解析归一化，$s$ 与 $1-s$ 对称，
故函数方程中不出现 $q^{-s}$ 的平移因子，这与第 10 章几何 Frobenius 的约定一致。
\par\vspace{0.15em}{\raggedleft 【来源：Hecke, \textit{\"Uber die Bestimmung Dirichletscher Reihen durch ihre Funktionalgleichung}, Math.\ Ann.\ 112 (1936), 664--699；Cogdell, \textit{Notes on $L$-functions for $\mathrm{GL}_n$}, ICTP Lecture Notes（A）】\par}
\end{annote}

关于 $\GL_n$ 的 $L$ 函数积分表示，还有另一批工作，出自代数 zeta 函数的经典研究。
这就是 Godement 与 Jacquet 发展的 $\GL_n$ 主 $L$ 函数理论 [15, 20]。
这条进路与本章所走的进路相关，但我们不打算介绍它。

这些 $L$ 函数的另一条进路，是经由 Eisenstein 级数的 Fourier 系数。
这条进路同样有经典渊源。
在自守表示的框架内，且在比 $\GL_n$ 更广的情形中，这一进路最初由 Langlands 铺陈 [29]，
随后由 Shahidi 最富于成效地推进。
Shahidi 在这方面的主要论文有 [35, 36, 37, 38, 39, 40, 41]。
特别地，他在 [38] 中证明两条进路给出相同的 $\GL_n$ 的 $L$ 函数。
本章不讨论这条进路。

\section{Fourier 展开}

本节中 $k$ 表整体域，$\A$ 表其阿代尔环，$\psi$ 表 $\A$ 上在 $k$ 上平凡的一个
连续加法特征。

我们从 $\GL_n(\A)$ 的一个尖点自守表示 $(\pi,V_\pi)$ 出发。
对我们而言，自守形式假定是光滑的（具有一致的中等增长），
但在阿基米德位置不必是 $K_\infty$-有限的。这对解析理论最为合用。
为简单起见，假定 $\pi$ 的中心特征 $\omega_\pi$ 是酉的。
于是 $V_\pi$ 是 $\GL_n(\A)$ 的某个不可约酉表示的光滑向量空间。
我们始终在上述意义下使用「尖点」一词：不可约酉尖点自守表示中的光滑向量。
（$\GL_n(\A)$ 的任一其他光滑尖点表示必形如 $\pi=\pi^\circ\otimes|\det|^t$，
其中 $\pi^\circ$ 酉、$t$ 实，故酉性假设确实不失一般性；它只是提供了一个方便的归一化。）
所谓 $\GL_n(\A)$ 上的尖点形式，指落在某个尖点表示中的函数。
所谓尖点函数，则指 $\GL_n(k)\backslash\GL_n(\A)$ 上的光滑函数 $\varphi$，
使对 $\GL_n$ 的标准抛物子群的每个单幂根 $U$，都有
$\int_{U(k)\backslash U(\A)}\varphi(ug)\,du\equiv0$。

本节的基本参考文献是 Piatetski-Shapiro [31, 32] 与 Shalika [42] 的论文。

\subsection{Fourier 展开式}

设 $f(\tau)$ 是上半平面 $\mathfrak{H}$ 上的全纯尖点形式，例如关于
$\mathrm{SL}_2(\mathbb{Z})$ 的。则 $f$ 在整平移下不变，$f(\tau+1)=f(\tau)$，
因而有形状如下的 Fourier 展开：
\[
f(\tau)=\sum_{n=1}^{\infty}a_ne^{2\pi in\tau}.
\]

若 $\varphi(g)$ 是 $\GL_2(\A)$ 上的光滑尖点形式，则平移对应于极大单幂子群
$N_2=\left\{n=\begin{pmatrix}1&x\\0&1\end{pmatrix}\right\}$，
且对 $n\in N_2(k)$ 有 $\varphi(ng)=\varphi(g)$。
于是，若 $\psi$ 是 $k\backslash\A$ 的任一连续特征，便可定义 $\psi$-Fourier 系数
或 $\psi$-Whittaker 函数
\[
W_{\varphi,\psi}(g)=\int_{k\backslash\A}\varphi\left(\begin{pmatrix}1&x\\0&1\end{pmatrix}g\right)\psi^{-1}(x)\,dx.
\]
相应的 Fourier 展开为
\[
\varphi(g)=\sum_{\psi}W_{\varphi,\psi}(g).
\]

（事实上，由 Abel Fourier 理论，作为 $x\in\A$ 的周期函数有
\[
\varphi\left(\begin{pmatrix}1&x\\0&1\end{pmatrix}g\right)=\sum_{\psi}W_{\varphi,\psi}(g)\psi(x),
\]
现在令 $x=0$ 即可。）

若固定 $k\backslash\A$ 的一个非平凡特征 $\psi$，则紧群 $k\backslash\A$ 的加法特征
通过映射 $\gamma\in k\mapsto\psi_\gamma$ 与 $k$ 同构，其中 $\psi_\gamma$ 由
$\psi_\gamma(x)=\psi(\gamma x)$ 定义。
一个初等计算表明，当 $\gamma\neq0$ 时
\[
W_{\varphi,\psi_\gamma}(g)=W_{\varphi,\psi}\left(\begin{pmatrix}\gamma&\\&1\end{pmatrix}g\right).
\]
若对固定的 $\psi$ 记 $W_\varphi=W_{\varphi,\psi}$，则 $\varphi$ 的 Fourier 展开成为
\[
\varphi(g)=W_{\varphi,\psi_0}(g)+\sum_{\gamma\in k^\times}W_\varphi\left(\begin{pmatrix}\gamma&\\&1\end{pmatrix}g\right).
\]

由于 $\varphi$ 是尖点的，
\[
W_{\varphi,\psi_0}(g)=\int_{k\backslash\A}\varphi\left(\begin{pmatrix}1&x\\0&1\end{pmatrix}g\right)dx\equiv0,
\]
而尖点形式 $\varphi$ 的 Fourier 展开就简化为
\[
\varphi(g)=\sum_{\gamma\in k^\times}W_\varphi\left(\begin{pmatrix}\gamma&\\&1\end{pmatrix}g\right).
\]

对 $\GL_n(\A)$ 上的尖点形式 $\varphi$，我们需要类似的展开。
平移仍对应于极大单幂子群
\[
N_n=\left\{n=\begin{pmatrix}
1&x_{1,2}&\cdots&*\\
&1&\ddots&\vdots\\
&&\ddots&x_{n-1,n}\\
0&&&1
\end{pmatrix}\right\},
\]
但此时它非交换。
这一困难由 Piatetski-Shapiro [31] 与 Shalika [42] 各自独立解决。
仍取上面固定的 $k\backslash\A$ 的非平凡连续特征 $\psi$，把它延拓为 $N_n$ 的特征，
令 $\psi(n)=\psi(x_{1,2}+\cdots+x_{n-1,n})$，并如下定义相应的 Fourier 系数
或 Whittaker 函数：
\[
W_\varphi(g)=W_{\varphi,\psi}(g)=\int_{N_n(k)\backslash N_n(\A)}\varphi(ng)\psi^{-1}(n)\,dn.
\]
由于 $\varphi$ 连续而积分区域紧，该积分绝对收敛，且在紧集上一致。
Fourier 展开具有下述形式。

\textbf{定理 1.1.}\quad 设 $\varphi\in V_\pi$ 是 $\GL_n(\A)$ 上的尖点形式，
$W_\varphi$ 是与之相伴的 $\psi$-Whittaker 函数。则
\[
\varphi(g)=\sum_{\gamma\in N_{n-1}(k)\backslash\GL_{n-1}(k)}
W_\varphi\left(\begin{pmatrix}\gamma&\\&1\end{pmatrix}g\right),
\]
且收敛绝对、在紧子集上一致。

此事的证明是一个归纳，用到 $\GL_n$ 的\emph{棒状子群}（mirabolic subgroup）$P_n$，
它似乎在 $\GL_n$ 自守形式的研究中无处不在。
抽象地说，$\GL_n$ 的棒状子群就是 $\GL_n$ 在 $k^n$ 上（任一）标准表示中
某个非零向量的稳定化子。
我们以 $P_n$ 记行向量 $e_n=(0,\ldots,0,1)\in k^n$ 的稳定化子，于是
\[
P_n=\left\{p=\begin{pmatrix}h&y\\0&1\end{pmatrix}\ \middle|\ h\in\GL_{n-1},\ y\in k^{n-1}\right\}
\cong\GL_{n-1}\ltimes Y_n,
\]
其中
\[
Y_n=\left\{y=\begin{pmatrix}I_{n-1}&y\\0&1\end{pmatrix}\ \middle|\ y\in k^{n-1}\right\}\cong k^{n-1}.
\]

仅由函数的限制，$\GL_n(\A)$ 上的尖点形式限制为 $P_n(\A)$ 上的光滑尖点函数，
且在 $P_n(k)$ 下仍左不变。
（$P_n(\A)$ 上在 $P_n(k)$ 下左不变的光滑函数 $\varphi$ 称为尖点的，
若对每个标准单幂子群 $U\subset P_n$ 都有
$\int_{U(k)\backslash U(\A)}\varphi(up)\,du\equiv0$。）
由于 $P_n\supset N_n$，可用与 $\GL_n(\A)$ 上相同的积分，
对 $P_n(\A)$ 上的尖点函数 $\varphi$ 定义相伴的 Whittaker 函数，即
\[
W_\varphi(p)=\int_{N_n(k)\backslash N_n(\A)}\varphi(np)\psi^{-1}(n)\,dn.
\]
对 $n$ 作归纳可以证明，对 $P_n(\A)$ 上的尖点函数 $\varphi$ 有
\[
\varphi(p)=\sum_{\gamma\in N_{n-1}(k)\backslash\GL_{n-1}(k)}
W_\varphi\left(\begin{pmatrix}\gamma&0\\0&1\end{pmatrix}p\right),
\]
且收敛绝对、在紧子集上一致。
\begin{annote}
这里正是 Gelfand--Kazhdan 与 Shalika 的贡献所在：$N_n$ 非交换，
通常的加法特征展开失效，替代物是 Whittaker 函数。
证明的骨架是对 $n$ 的归纳，而归纳的载体是棒状子群 $P_n$：
它的单幂根 $N_n$ 足以定义 Whittaker 函数，而商 $P_n/N_n\cong\GL_{n-1}$
又让归纳假设可以调用。
棒状子群之所以「无处不在」，是因为 $P_n\backslash\GL_n\cong k^n-\{0\}$，
而 $k^n-\{0\}$ 上的轨道分析恰好把 $N_n$ 的积分与 $e_n$ 的稳定化子对接起来。
\par\vspace{0.15em}{\raggedleft 【来源：Shalika, \textit{The multiplicity one theorem for $\mathrm{GL}(n)$}, Ann.\ of Math.\ 100 (1974), 171--193；Piatetski-Shapiro, \textit{Euler subgroups}, 1971（A）】\par}
\end{annote}

要从这个展开得到 $\GL_n$ 上的 Fourier 展开，设 $\varphi$ 是 $\GL_n(\A)$ 上的尖点形式。
对 $g\in\Omega$（$\Omega$ 为紧子集），函数 $\varphi_g(p)=\varphi(pg)$
构成 $P_n(\A)$ 上的一族紧尖点函数。于是
\[
\varphi_g(1)=\sum_{\gamma\in N_{n-1}(k)\backslash\GL_{n-1}(k)}
W_{\varphi_g}\left(\begin{pmatrix}\gamma&0\\0&1\end{pmatrix}\right),
\]
收敛绝对且一致。因此
\[
\varphi(g)=\sum_{\gamma\in N_{n-1}(k)\backslash\GL_{n-1}(k)}
W_\varphi\left(\begin{pmatrix}\gamma&0\\0&1\end{pmatrix}g\right),
\]
同样绝对收敛，且对 $g\in\Omega$ 一致。

\subsection{Whittaker 模型与重数一}

现在考察尖点形式 $\varphi$ 的 Fourier 展开中出现的函数 $W_\varphi$。
它们都是 $\GL_n(\A)$ 上的光滑函数 $W(g)$，且对 $n\in N_n(\A)$ 满足
$W(ng)=\psi(n)W(g)$。
若令 $\mathcal{W}(\pi,\psi)=\{W_\varphi\mid\varphi\in V_\pi\}$，
则 $\GL_n(\A)$ 通过右平移作用于此空间，而映射 $\varphi\mapsto W_\varphi$
把 $V_\pi$ 与 $\mathcal{W}(\pi,\psi)$ 缠结起来。
$\mathcal{W}(\pi,\psi)$ 称为 $\pi$ 的 \emph{Whittaker 模型}。

表示的 Whittaker 模型这一概念，在局部域上完全说得通。
设 $k_v$ 为局部域（例如 $k$ 的一个完备化），$(\pi_v,V_{\pi_v})$ 为
$\GL_n(k_v)$ 的不可约容许光滑表示。
固定 $k_v$ 的一个非平凡连续加法特征 $\psi_v$。
以 $\mathcal{W}(\psi_v)$ 记 $\GL_n(k_v)$ 上所有满足
$W(ng)=\psi_v(n)W(g)$（对一切 $n\in N_n(k_v)$）的光滑函数 $W(g)$ 之空间，
即 $\GL_n(k_v)$ 上关于 $\psi_v$ 的所有光滑 Whittaker 函数之空间。
这也就是光滑诱导表示 $\operatorname{Ind}_{N_n}^{\GL_n}(\psi_v)$ 的空间。
$\GL_n(k_v)$ 通过右平移作用于其上。
若有一个非平凡的连续 intertwining $V_{\pi_v}\to\mathcal{W}(\psi_v)$，
则记其像为 $\mathcal{W}(\pi_v,\psi_v)$，并称之为 $\pi_v$ 的 Whittaker 模型。

表示 $(\pi_v,V_{\pi_v})$ 的 Whittaker 模型，等价于 $V_{\pi_v}$ 上的连续
Whittaker 泛函，即满足
$\Lambda_v(\pi_v(n)\xi_v)=\psi_v(n)\Lambda_v(\xi_v)$（对一切 $n\in N_n(k_v)$）
的连续泛函 $\Lambda_v$。
由模型得泛函，只需令 $\Lambda_v(\xi_v)=W_{\xi_v}(e)$；
由泛函得模型，只需令 $W_{\xi_v}(g)=\Lambda_v(\pi_v(g)\xi_v)$。
这是 Frobenius 互反的一种形式，在此即同构
$\Hom_{N_n}(V_{\pi_v},\mathbb{C}_{\psi_v})\cong\Hom_{\GL_n}(V_{\pi_v},\operatorname{Ind}_{N_n}^{\GL_n}(\psi_v))$，
其构造如上。

关于 Whittaker 泛函与模型的存在唯一性，基本定理如下。

\textbf{定理 1.2.}\quad 设 $(\pi_v,V_{\pi_v})$ 是 $\GL_n(k_v)$ 的不可约容许光滑表示，
$\psi_v$ 是 $k_v$ 的非平凡连续加法特征。则 $V_{\pi_v}$ 上连续 $\psi_v$-Whittaker
泛函之空间至多一维。也就是说，Whittaker 模型若存在，则唯一。

这一结果最先由 Gelfand 与 Kazhdan [14] 对非阿基米德域证明，
其后由 Shalika [42] 推广到阿基米德局部域。

$\GL_n(k_v)$ 的不可约容许光滑表示 $(\pi_v,V_{\pi_v})$ 若具有 Whittaker 模型，
则称为\emph{一般的}（generic）或\emph{非退化的}。
Gelfand 与 Kazhdan 还证明：$\pi_v$ 一般当且仅当其反步表示 $\widetilde{\pi}_v$ 一般；
事实上 $\widetilde{\pi}\cong\pi^\iota$，其中 $\iota$ 是外自同构 $g^\iota={}^tg^{-1}$，
而此时 $\widetilde{\pi}_v$ 的 Whittaker 模型可取为
\[
\mathcal{W}(\widetilde{\pi}_v,\psi_v^{-1})=\left\{\widetilde{W}(g)=W(w_n{}^tg^{-1})\ \middle|\ W\in\mathcal{W}(\pi,\psi_v)\right\}.
\]

由 Whittaker 模型的局部唯一性，可以推出整体唯一性。
若 $(\pi,V_\pi)$ 是 $\GL_n(\A)$ 的不可约光滑容许表示，
则 $\pi$ 分解为局部表示的受限张量积 $\pi\cong\otimes'\pi_v$，取遍 $k$ 的所有位置 [10, 13]。
于是对每个局部分量都有连续嵌入 $V_{\pi_v}\hookrightarrow V_\pi$。
因此 $V_\pi$ 上任一 Whittaker 泛函 $\Lambda$ 决定各 $V_{\pi_v}$ 上的一族局部
Whittaker 泛函 $\Lambda_v$，反之亦然，使得 $\Lambda=\otimes'\Lambda_v$。
故整体唯一性可由局部唯一性推出。
此外，一旦固定 $V_\pi$ 与 $\otimes'V_{\pi_v}$ 的同构，并用 $\Lambda$ 及相应的
$\Lambda_v$ 族定义整体与局部 Whittaker 函数，便有整体 Whittaker 函数的分解
\[
W_\xi(g)=\prod_v W_{\xi_v}(g_v),
\]
其中 $\xi\in V_\pi$ 可分解，即 $\xi=\otimes'\xi_v$ 对应一个纯张量。
我们将会看到，这一分解是 Whittaker 模型唯一性的直接推论，
在 $\GL_n$ 的 Euler 积分理论中起极为重要的作用。
\begin{annote}
Whittaker 模型的唯一性之所以是整章的枢纽，是因为它把「整体」与「局部」的
$L$ 函数对接起来：整体积分一旦展开成 Whittaker 函数的积分，唯一性就保证
被积函数是纯张量上的乘积，积分遂自动分解为局部积分之积。
反过来，若没有唯一性，整体积分仍可定义，却无法保证它是 Euler 乘积。
另一处常见混淆是反步表示上的 $\psi^{-1}$：正文把 $\mathcal{W}(\widetilde{\pi}_v,\psi_v^{-1})$
写成 $W(w_n{}^tg^{-1})$ 的集合，这里的 $w_n$ 是长 Weyl 元，
而 $g\mapsto w_n{}^tg^{-1}$ 正是把左 $\psi$-等变换成右 $\psi^{-1}$-等变的那一步。
\par\vspace{0.15em}{\raggedleft 【来源：Gelfand--Kazhdan, \textit{Representations of $\mathrm{GL}(n,K)$ where $K$ is a local field}, 1971；Shalika, Ann.\ of Math.\ 100 (1974)（A）】\par}
\end{annote}

现在来看这对 $\GL_n(\A)$ 的尖点表示 $(\pi,V_\pi)$ 意味着什么。
我们已经看到，对任一光滑尖点形式 $\varphi\in V_\pi$，有 Fourier 展开
\[
\varphi(g)=\sum_{\gamma\in N_{n-1}(k)\backslash\GL_{n-1}(k)}
W_\varphi\left(\begin{pmatrix}\gamma&\\&1\end{pmatrix}g\right).
\]
由此可以断定 $\mathcal{W}(\pi,\psi)\neq0$，且 $\pi$ 是（整体）一般的，
其 Whittaker 泛函为
\[
\Lambda(\varphi)=W_\varphi(e)=\int\varphi(ng)\psi^{-1}(n)\,dn.
\]
于是 $\varphi$ 完全由相伴的 Whittaker 函数 $W_\varphi$ 决定。
由整体 Whittaker 模型的唯一性，可以导出 Piatetski-Shapiro [32] 与
Shalika [42] 的重数一定理。

\textbf{重数一.}\quad 设 $(\pi,V_\pi)$ 是 $\GL_n(\A)$ 的不可约光滑容许表示。
则 $\pi$ 在 $\GL_n(\A)$ 上尖点形式空间中的重数至多为一。

\section{Euler 积分表示}

仍设 $f(\tau)$ 是全模群对应的 $\mathfrak{H}$ 上权为 $k$ 的全纯尖点形式，
其 Fourier 展开为
\[
f(\tau)=\sum a_ne^{2\pi in\tau}.
\]
Hecke [18] 于是给 $f$ 配上 $L$ 函数
\[
L(s,f)=\sum a_nn^{-s},
\]
并通过把 $L(s,f)$ 写成 $f$ 的 Mellin 变换，分析其解析性质，
即延拓、增长阶与函数方程：
\[
\Lambda(s,f)=(2\pi)^{-s}\Gamma(s)L(s,f)=\int_0^\infty f(iy)y^s\,d^\times y.
\]
对 $f(\tau)$ 在变换 $\tau\mapsto-1/\tau$ 下的模变换律作一应用，便得到函数方程
\[
\Lambda(s,f)=(-1)^{k/2}\Lambda(k-s,f).
\]
此外，若 $f$ 是所有 Hecke 算子的本征函数，则 $L(s,f)$ 有 Euler 乘积展开
\[
L(s,f)=\prod_p(1-a_pp^{-s}+p^{k-1-2s})^{-1}.
\]

对 $\GL_n(\A)$ 的尖点自守表示 $(\pi,V_\pi)$，有类似的理论。
为了应用 Langlands 猜想，并经由逆定理应用函子性，我们不仅需要标准的
$L$ 函数 $L(s,\pi)$，还需要扭 $L$ 函数 $L(s,\pi\times\pi')$，
其中 $(\pi',V_{\pi'})$ 是 $\GL_m(\A)$ 的尖点自守表示，$m<n$。

本节的基本参考文献是 Jacquet--Langlands [21]、Jacquet、Piatetski-Shapiro 与
Shalika [22]，以及 Jacquet 与 Shalika [25]。

\subsection{$\GL_2$ 的积分表示}

先考虑 $\GL_2(\A)$ 的尖点自守表示 $(\pi,V_\pi)$ 的 $L$ 函数，
扭以一个 id\`ele 类特征 $\chi$；这等价于扭以 $\GL_1(\A)$ 的一个（尖点）自守表示，
如 Jacquet--Langlands [21] 中那样。
遵循 Jacquet 与 Langlands（他们又遵循 Hecke），对每个 $\varphi\in V_\pi$ 考虑积分
\[
I(s;\varphi,\chi)=\int_{k^\times\backslash\A^\times}
\varphi\begin{pmatrix}a&\\&1\end{pmatrix}\chi(a)|a|^{s-1/2}\,d^\times a.
\]
由于 $\GL_2(\A)$ 上的尖点形式限制到 $\A^\times$ 上（如积分中所示）速降，
该积分对一切 $s$ 绝对收敛，且在 $\operatorname{Re}(s)$ 的任一区间上一致。
故 $I(s;\varphi,\chi)$ 是 $s$ 的整函数，在任一竖直带 $a\le\operatorname{Re}(s)\le b$ 上有界。
此外，若令 $\widetilde{\varphi}(g)=\varphi({}^tg^{-1})=\varphi(w_n{}^tg^{-1})$，
则 $\widetilde{\varphi}\in V_{\widetilde{\pi}}$；而积分中作变量替换 $a\mapsto a^{-1}$，
便知每个积分满足形状如下的函数方程
\[
I(s;\varphi,\chi)=I(1-s;\widetilde{\varphi},\chi^{-1}).
\]
可见这些积分各自具有相当好的解析性质。
\begin{annote}
这里的 $|a|^{s-1/2}$ 不是随意选的归一化。
若用 $|a|^{s}$，局部积分对应 $L(s+1/2,\pi\times\chi)$ 型的平移；
取 $s-1/2$ 是为了让 $\GL_2$ 的中心特征与 $\GL_1$ 的特征在函数方程中对称出现，
使 $s\mapsto1-s$ 恰好对调两个因子。
Tate 论文的原始做法是对 $\GL_1$ 用 $|a|^s$，本节的 $1/2$ 平移是 $\GL_2$
相对于 $\GL_1$ 的「半权」修正。
\par\vspace{0.15em}{\raggedleft 【来源：Tate, \textit{Fourier analysis in number fields and Hecke's zeta-functions}, 1950；Jacquet--Langlands, \textit{Automorphic Forms on $\mathrm{GL}(2)$}, Springer LNM 114, 1970（A）】\par}
\end{annote}

若把 $\varphi$ 换成它在第 1 节的 Fourier 展开再展开积分，便得
\[
\begin{aligned}
I(s;\varphi,\chi)&=\int_{k^\times\backslash\A^\times}
\sum_{\gamma\in k^\times}W_\varphi\left(\begin{pmatrix}\gamma a&\\&1\end{pmatrix}\right)
\chi(a)|a|^{s-1/2}\,d^\times a\\
&=\int_{\A^\times}W_\varphi\left(\begin{pmatrix}a&\\&1\end{pmatrix}\right)
\chi(a)|a|^{s-1/2}\,d^\times a,
\end{aligned}
\]
这里用到函数 $\chi(a)|a|^{s-1/2}$ 在 $k^\times$ 下不变。
由 Whittaker 函数的标准规范估计（gauge estimate）[22]，展开后该式对
$\operatorname{Re}(s)\gg0$ 收敛。
如第 1 节所见，若 $W_\varphi\in\mathcal{W}(\pi,\psi)$ 对应于可分解向量
$\varphi\in V_\pi\cong\otimes'V_{\pi_v}$，则 Whittaker 函数分解为局部
Whittaker 函数之积
\[
W_\varphi(g)=\prod_v W_{\varphi_v}(g_v).
\]
由于特征 $\chi$ 与阿代尔绝对值都分解为局部分量，积分区域 $\A^\times$ 也分解，
故整体积分自然地分解为局部积分之积
\[
\int_{\A^\times}W_\varphi\left(\begin{pmatrix}a&\\&1\end{pmatrix}\right)\chi(a)|a|^{s-1/2}\,d^\times a
=\prod_v\int_{k_v^\times}W_{\varphi_v}\left(\begin{pmatrix}a_v&\\&1\end{pmatrix}\right)
\chi_v(a_v)|a_v|^{s-1/2}\,d^\times a_v,
\]
其中无穷乘积对 $\operatorname{Re}(s)\gg0$ 仍收敛；或写作
\[
I(s;\varphi,\chi)=\prod_v\Psi_v(s;W_{\varphi_v},\chi_v),
\]
局部积分按显然的方式定义为
\[
\Psi_v(s;W_{\varphi_v},\chi_v)=\int_{k_v^\times}
W_{\varphi_v}\left(\begin{pmatrix}a_v&\\&1\end{pmatrix}\right)
\chi_v(a_v)|a_v|^{s-1/2}\,d^\times a_v.
\]
于是每个整体积分都是 Euler 乘积。
这样，对 $\pi$ 与 $\chi$ 便配上了一族具有良好解析性质的整体 Euler 积分，
同时对每个位置 $v$ 配上了一族对 $\operatorname{Re}(s)\gg0$ 收敛的局部积分。

\subsection{$m<n$ 时 $\GL_n\times\GL_m$ 的积分表示}

现在设 $(\pi,V_\pi)$ 是 $\GL_n(\A)$ 的尖点表示，$(\pi',V_{\pi'})$ 是
$\GL_m(\A)$ 的尖点表示，$m<n$。取 $\varphi\in V_\pi$，$\varphi'\in V_{\pi'}$。
乍看之下，$\GL_2$ 配 $\GL_1$ 扭的积分，其自然类比应是
\[
\int_{\GL_m(k)\backslash\GL_m(\A)}
\varphi\begin{pmatrix}h&\\&I_{n-m}\end{pmatrix}\varphi'(h)|\det(h)|^{s-(n-m)/2}\,dh.
\]
这族积分确如前具有一切良好解析性质（有限阶整函数，满足函数方程），
但除 $m=n-1$ 外都不是 Euler 乘积；$m=n-1$ 的情形与 $\GL_2$ 情形完全一样。
问题在于，把形式 $\varphi$ 限制到 $\GL_m$ 上过于粗暴，
代入 $\varphi$ 的 Fourier 展开后无法得到漂亮的展开。
我们转而引入投影算子，把 $\GL_n(\A)$ 上的尖点形式投到 $P_{m+1}(\A)$ 上的尖点函数，
这些算子由定义 Whittaker 函数时所用单幂积分的一部分给出。

在 $\GL_n$ 中，以 $Y_{n,m}$ 记对应于分拆 $(m+1,1,\ldots,1)$ 的标准抛物子群的单幂根。
若 $\psi$ 是 $k\backslash\A$ 的标准加法特征，则 $\psi$ 定义了 $Y_{n,m}(\A)$ 的一个
在 $Y_{n,m}(k)$ 上平凡的特征，因为 $Y_{n,m}\subset N_n$。
群 $Y_{n,m}$ 被 $\GL_{m+1}\subset\GL_n$ 归一化，
而棒状子群 $P_{m+1}\subset\GL_{m+1}$ 正是 $\GL_{m+1}$ 中特征 $\psi$ 的稳定化子。

若 $\varphi(g)$ 是 $\GL_n(\A)$ 上的尖点形式，则如下定义投影算子
$\mathbb{P}^n_m$，它把 $\GL_n(\A)$ 上的尖点形式送到 $P_{m+1}(\A)$ 上的尖点函数：
\[
\mathbb{P}^n_m\varphi(p)=|\det(p)|^{-\left(\frac{n-m-1}{2}\right)}
\int_{Y_{n,m}(k)\backslash Y_{n,m}(\A)}
\varphi\left(y\begin{pmatrix}p&\\&I_{n-m}\end{pmatrix}\right)\psi^{-1}(y)\,dy,
\]
其中 $p\in P_{m+1}(\A)$。
由于积分区域紧，积分绝对收敛。容易验证 $\mathbb{P}^n_m\varphi(p)$ 确为
$P_{m+1}(\A)$ 上的尖点函数。
由第 1 节知，$P_{m+1}(\A)$ 上的尖点函数有对 $N_m(k)\backslash\GL_m(k)$ 求和的
Fourier 展开。把此展开用于投影后的 $\GL_n(\A)$ 尖点形式，
便知对 $h\in\GL_m(\A)$，$\mathbb{P}^n_m\varphi\begin{pmatrix}h&\\&1\end{pmatrix}$ 有展开
\[
\mathbb{P}^n_m\varphi\begin{pmatrix}h&\\&1\end{pmatrix}
=|\det(h)|^{-\left(\frac{n-m-1}{2}\right)}
\sum_{\gamma\in N_m(k)\backslash\GL_m(k)}
W_\varphi\left(\begin{pmatrix}\gamma&0\\0&I_{n-m}\end{pmatrix}
\begin{pmatrix}h&\\&I_{n-m}\end{pmatrix}\right),
\]
且收敛绝对、在紧子集上一致。
\begin{annote}
投影算子 $\mathbb{P}^n_m$ 是本章技术上最关键、也最容易被读者跳过的一步。
它的作用是把 $\varphi$ 先沿 $Y_{n,m}$ 作一次单幂积分，
从而把 $\varphi$ 在 $\GL_n$ 上的 Fourier 展开「降」到 $P_{m+1}$ 上；
这样再做 $\GL_m$ 上的积分时，展开才可能化成单个 Whittaker 函数。
归一化因子 $|\det(p)|^{-(n-m-1)/2}$ 保证 $\mathbb{P}^n_m$ 与 $P_{m+1}$ 上的
模特征相容，缺了它，后面 $\Psi_v$ 的指数 $s-(n-m)/2$ 就对不上。
注意 $Y_{n,m}$ 与 $N_n$ 的关系：前者是后者的子群，故 $\psi$ 的限制仍有意义。
\par\vspace{0.15em}{\raggedleft 【来源：Jacquet--Piatetski-Shapiro--Shalika, \textit{Rankin--Selberg convolutions}, Amer.\ J.\ Math.\ 105 (1983), 367--464（A）】\par}
\end{annote}

现在我们已具备写下 $\GL_n$ 上尖点形式被 $\GL_m$ 上自守形式扭的 Euler 积分族的条件，
其中 $m<n$。设 $\varphi\in V_\pi$ 是 $\GL_n(\A)$ 上的尖点形式，
$\varphi'\in V_{\pi'}$ 是 $\GL_m(\A)$ 上的尖点形式。
（其实可以取 $\varphi'$ 为 $\GL_m(\A)$ 上任一自守形式。）
考虑积分
\[
I(s;\varphi,\varphi')=\int_{\GL_m(k)\backslash\GL_m(\A)}
\mathbb{P}^n_m\varphi\begin{pmatrix}h&0\\0&1\end{pmatrix}
\varphi'(h)|\det(h)|^{s-1/2}\,dh.
\]
由于尖点形式速降，积分 $I(s;\varphi,\varphi')$ 对复参数 $s$ 的一切值绝对收敛，
且在紧子集上一致。故如前，它是整函数，并在任一竖直带上有界。

下面考察这些积分的 Euler 性质。先用 $\mathbb{P}^n_m\varphi$ 的 Fourier 展开替换之，得
\[
\begin{aligned}
I(s;\varphi,\varphi')
&=\int_{\GL_m(k)\backslash\GL_m(\A)}
\sum_{\gamma\in N_m(k)\backslash\GL_m(k)}
W_\varphi\left(\begin{pmatrix}\gamma h&0\\0&I_{n-m}\end{pmatrix}\right)
\varphi'(h)|\det(h)|^{s-(n-m)/2}\,dh.
\end{aligned}
\]
由于 $\varphi'(h)$ 在 $\GL_m(\A)$ 上自守，且对 $\gamma\in\GL_m(k)$ 有
$|\det(\gamma)|=1$，故当 $\operatorname{Re}(s)\gg0$ 时可交换求和与积分次序，
再合并同类项，得
\[
I(s;\varphi,\varphi')=\int_{N_m(k)\backslash\GL_m(\A)}
W_\varphi\begin{pmatrix}h&0\\0&I_{n-m}\end{pmatrix}
\varphi'(h)|\det(h)|^{s-(n-m)/2}\,dh.
\]
由 [22, 第 13 节] 的规范估计，该积分对 $\operatorname{Re}(s)\gg0$ 绝对收敛，
交换次序由此得到正当化。
现在先对 $N_m(k)\backslash N_m(\A)$ 积分。
注意对 $n\in N_m(\A)\subset N_n(\A)$ 有 $W_\varphi(ng)=\psi(n)W_\varphi(g)$，故得
\[
\begin{aligned}
I(s;\varphi,\varphi')
&=\int_{N_m(\A)\backslash\GL_m(\A)}
W_\varphi\begin{pmatrix}h&0\\0&I_{n-m}\end{pmatrix}
W'_{\varphi'}(h)|\det(h)|^{s-(n-m)/2}\,dh\\
&=\Psi(s;W_\varphi,W'_{\varphi'}),
\end{aligned}
\]
其中 $W'_{\varphi'}(h)$ 是与 $\varphi'$ 相伴的 $\GL_m(\A)$ 上 $\psi^{-1}$-Whittaker 函数，即
\[
W'_{\varphi'}(h)=\int_{N_m(k)\backslash N_m(\A)}\varphi'(nh)\psi(n)\,dn,
\]
而 $\operatorname{Re}(s)\gg0$ 时的绝对收敛仍然保持。

到此为止，积分之为 Euler 乘积，已是 $\GL_n$ Whittaker 模型唯一性的推论。
取 $\GL_n(\A)$ 的尖点表示 $\pi$ 中的光滑尖点形式 $\varphi$，
并额外假定 $\varphi$ 可分解：在 $\pi$ 分解为局部表示受限张量积
$\pi=\otimes'\pi_v$ 的分解中，$\varphi=\otimes\varphi_v$ 是纯张量。
则如前所见，可选取局部 Whittaker 模型使
$W_\varphi(g)=\prod_v W_{\varphi_v}(g_v)$。
对可分解的 $\varphi'$ 同理有分解 $W'_{\varphi'}(h)=\prod_v W'_{\varphi'_v}(h_v)$。
把这些分解代入积分式，并注意积分区域也分解为
$N_m(\A)\backslash\GL_m(\A)=\prod_v N_m(k_v)\backslash\GL_m(k_v)$，
即知积分分解为局部积分之积
\[
I(s;\varphi,\varphi')=\Psi(s;W_\varphi,W'_{\varphi'})
=\prod_v\Psi_v(s;W_{\varphi_v},W'_{\varphi'_v}),
\]
其中局部积分为
\[
\Psi_v(s;W_{\varphi_v},W'_{\varphi'_v})
=\int_{N_m(k_v)\backslash\GL_m(k_v)}
W_{\varphi_v}\begin{pmatrix}h_v&0\\0&I_{n-m}\end{pmatrix}
W'_{\varphi'_v}(h_v)|\det(h_v)|_v^{s-(n-m)/2}\,dh_v.
\]
由 [22, 命题 2.3.6] 的规范估计，各局部积分对 $\operatorname{Re}(s)\gg0$ 收敛。
至此我们已构造出一族 Euler 积分。

现在回到函数方程问题。与 $\GL_2$ 情形一样，函数方程本质上是外自同构
$g\mapsto\iota(g)=g^\iota={}^tg^{-1}$ 存在的结果。
若用此自同构作用于自守形式，令
$\widetilde{\varphi}(g)=\varphi(g^\iota)=\varphi(w_n{}^tg^{-1})$，
并令 $\widetilde{\mathbb{P}}^n_m=\iota\circ\mathbb{P}^n_m\circ\iota$，
则积分自然地满足函数方程
\[
I(s;\varphi,\varphi')=\widetilde{I}(1-s;\widetilde{\varphi},\widetilde{\varphi}'),
\]
其中
\[
\widetilde{I}(s;\varphi,\varphi')=\int_{\GL_m(k)\backslash\GL_m(\A)}
\widetilde{\mathbb{P}}^n_m\varphi\begin{pmatrix}h&\\&1\end{pmatrix}
\varphi'(h)|\det(h)|^{s-1/2}\,dh.
\]

我们已建立了下述结果。

\textbf{定理 2.1.}\quad 设 $\varphi\in V_\pi$ 是 $\GL_n(\A)$ 上的尖点形式，
$\varphi'\in V_{\pi'}$ 是 $\GL_m(\A)$ 上的尖点形式，$m<n$。
则积分族 $I(s;\varphi,\varphi')$ 定义 $s$ 的整函数，在竖直带上有界，
并满足函数方程
\[
I(s;\varphi,\varphi')=\widetilde{I}(1-s;\widetilde{\varphi},\widetilde{\varphi}').
\]
此外这些积分是 Euler 乘积；若 $\varphi$ 与 $\varphi'$ 可分解，则
\[
I(s;\varphi,\varphi')=\prod_v\Psi_v(s;W_{\varphi_v},W'_{\varphi'_v}),
\]
收敛绝对且对 $\operatorname{Re}(s)\gg0$ 一致。

函数方程右端出现的积分同样是 Euler 乘积。展开定义，先得到
\[
\widetilde{I}(1-s;\widetilde{\varphi},\widetilde{\varphi}')
=\widetilde{\Psi}(1-s;\rho(w_{n,m})\widetilde{W}_\varphi,\widetilde{W}'_{\varphi'}),
\]
其中展开后的整体积分为
\[
\widetilde{\Psi}(s;W,W')=\iint W\left(\begin{pmatrix}h\\x&I_{n-m-1}\\&&1\end{pmatrix}\right)
dx\,W'(h)|\det(h)|^{s-(n-m)/2}\,dh,
\]
$h$ 的积分取在 $N_m(\A)\backslash\GL_m(\A)$ 上，
$x$ 的积分取在 $M_{n-m-1,m}(\A)$ 上，即 $(n-m-1)\times m$ 矩阵之空间，
$\rho$ 表右平移，$w_{n,m}$ 是 Weyl 元
\[
w_{n,m}=\begin{pmatrix}I_m&\\&w_{n-m}\end{pmatrix},
\qquad
w_{n-m}=\begin{pmatrix}&&1\\&\cdot&\\1&&\end{pmatrix},
\]
$w_{n-m}$ 是 $\GL_{n-m}$ 中的标准长 Weyl 元。
又对 $W\in\mathcal{W}(\pi,\psi)$，置
$\widetilde{W}(g)=W(w_n{}^tg^{-1})\in\mathcal{W}(\pi,\psi^{-1})$。
多出的单幂积分是 $\widetilde{\mathbb{P}}^n_m$ 的残余。
如前，$\widetilde{\Psi}(s;W,W')$ 对 $\operatorname{Re}(s)\gg0$ 绝对收敛。
若 $\varphi$ 与 $\varphi'$ 如前可分解，则这些积分 $\widetilde{\Psi}(s;W_\varphi,W'_{\varphi'})$
同样分解。于是有
\[
\widetilde{\Psi}(s;W_\varphi,W'_{\varphi'})
=\prod_v\widetilde{\Psi}_v(s;W_{\varphi_v},W'_{\varphi'_v}),
\]
其中
\[
\widetilde{\Psi}_v(s;W_v,W'_v)=\iint W_v\left(\begin{pmatrix}h_v\\x_v&I_{n-m-1}\\&&1\end{pmatrix}\right)
dx_v\,W'_v(h_v)|\det(h_v)|^{s-(n-m)/2}\,dh_v,
\]
此时 $h_v$ 的积分取在 $N_m(k_v)\backslash\GL_m(k_v)$ 上，
$x_v$ 的积分取在矩阵空间 $M_{n-m-1,m}(k_v)$ 上。
回到函数方程，即知右端是 Euler 乘积，并分解为
\[
\widetilde{I}(1-s;\widetilde{\varphi},\widetilde{\varphi}')
=\widetilde{\Psi}(1-s;\rho(w_{n,m})\widetilde{W}_\varphi,\widetilde{W}'_{\varphi'})
=\prod_v\widetilde{\Psi}_v(1-s;\rho(w_{n,m})\widetilde{W}_{\varphi_v},\widetilde{W}'_{\varphi'_v}).
\]

\subsection{$\GL_n\times\GL_n$ 的积分表示}

$\GL_n\times\GL_n$ 的 $L$ 函数积分表示的范式并非来自 Hecke，
而是来自 Rankin [33] 与 Selberg [34] 的经典论文。
这些工作最先由 Jacquet 在 $\GL_2\times\GL_2$ 情形纳入自守表示的框架 [20]，
随后由 Jacquet 与 Shalika 处理一般情形 [25]。

设 $(\pi,V_\pi)$ 与 $(\pi',V_{\pi'})$ 是 $\GL_n(\A)$ 的两个尖点表示。
取两个尖点形式 $\varphi\in V_\pi$，$\varphi'\in V_{\pi'}$。
上述构造的类比将只是
\[
\int_{\GL_n(k)\backslash\GL_n(\A)}\varphi(g)\varphi'(g)|\det(g)|^s\,dg.
\]
这个积分本质上是 $\varphi$ 与 $\varphi'$ 的 $L^2$ 内积，
不适合用来定义 $L$ 函数，尽管它会作为我们积分在极点处的留数出现。
我们改循 Rankin 与 Selberg 的做法，采用一个涉及第三个函数的积分表示，
即 $\GL_n(\A)$ 上的一个 Eisenstein 级数。
这族 Eisenstein 级数仍用棒状子群来构造。

为构造 Eisenstein 级数，回到观察 $P_n\backslash\GL_n\cong k^n-\{0\}$。
以 $\mathcal{S}(\A^n)$ 记 $\A^n$ 上的 Schwartz--Bruhat 函数空间，
则每个 $\Phi\in\mathcal{S}$ 通过 $g\mapsto\Phi((0,\ldots,0,1)g)=\Phi(e_ng)$
定义了 $\GL_n(\A)$ 上一个在 $P_n(\A)$ 下左不变的光滑函数。
设 $\eta$ 是酉 id\`ele 类特征。（在我们的应用中，$\eta$ 将由 $\pi$ 与 $\pi'$
的中心特征决定。）考虑函数
\[
F(g,\Phi;s,\eta)=|\det(g)|^s\int_{\A^\times}\Phi(ae_ng)|a|^{ns}\eta(a)\,d^\times a.
\]
若令 $P'_n=Z_nP_n$ 为 $\GL_n$ 中对应于分拆 $(n-1,1)$ 的抛物子群，
并把 $\eta$ 延拓为 $P'_n$ 的特征，令 $\eta(p')=\eta(d)$，
其中 $p'=\begin{pmatrix}h&\\&d\end{pmatrix}\in P'_n(\A)$，$h\in\GL_{n-1}(\A)$，$d\in\A^\times$，
则 $F(g,\Phi;s,\eta)$ 是正规化诱导表示
$\operatorname{Ind}_{P'_n(\A)}^{\GL_n(\A)}(\delta_{P'_n}^{s-1/2}\eta)$ 的光滑截面。
由于 $P'_n(\A)$ 的诱导特征 $\delta_{P'_n}^{s-1/2}\eta$ 在 $P'_n(k)$ 下不变，
便可由这族截面构造 Eisenstein 级数
\[
E(g,\Phi;s,\eta)=\sum_{\gamma\in P'_n(k)\backslash\GL_n(k)}F(\gamma g,\Phi;s,\eta).
\]
此式对 $\operatorname{Re}(s)>1$ 绝对收敛 [25]。

若把和式中的 $F$ 换成其定义并展开，可把这 Eisenstein 级数改写为
\[
E(g,\Phi;s,\eta)=|\det(g)|^s\int_{k^\times\backslash\A^\times}
\Theta'_\Phi(a,g)|a|^{ns}\eta(a)\,d^\times a.
\]
第二个表达式本质上把 Eisenstein 级数表为 Theta 级数的 Mellin 变换，
该 Theta 级数为
\[
\Theta_\Phi(a,g)=\sum_{\xi\in k^n}\Phi(a\xi g),
\]
其中上文已置
\[
\Theta'_\Phi(a,g)=\sum_{\xi\in k^n-\{0\}}\Phi(a\xi g)=\Theta_\Phi(a,g)-\Phi(0).
\]
这使我们能由 $\Theta_\Phi$ 的 Poisson 求和公式得到 Eisenstein 级数的解析性质。
Poisson 求和给出
\[
\begin{aligned}
E(g,\Phi,s,\eta)&=|\det(g)|^s\int_{|a|\ge1}\Theta'_\Phi(a,g)|a|^{ns}\eta(a)\,d^\times a\\
&\quad+|\det(g)|^{s-1}\int_{|a|\ge1}\Theta'_{\widehat{\Phi}}(a,g^\iota)
|a|^{n(1-s)}\eta^{-1}(a)\,d^\times a+\delta(s),
\end{aligned}
\]
其中 $\mathcal{S}(\A^n)$ 上的 Fourier 变换 $\widehat{\Phi}$ 定义为
\[
\widehat{\Phi}(x)=\int_{\A^\times}\Phi(y)\psi(y{}^tx)\,dy,
\]
而
\[
\delta(s)=\begin{cases}
0,&\text{若 }\eta\text{ 分歧；}\\[2pt]
-c\Phi(0)\dfrac{|\det(g)|^s}{s+i\sigma}
+c\widehat{\Phi}(0)\dfrac{|\det(g)|^{s-1}}{s-1+i\sigma},&
\text{若 }\eta(a)=|a|^{in\sigma},\ \sigma\in\mathbb{R},
\end{cases}
\]
$c$ 为非零常数。

由此容易推出 Eisenstein 级数的基本性质 [25, 第 4 节]。
Eisenstein 级数 $E(g,\Phi;s,\eta)$ 可亚纯延拓到整个 $\mathbb{C}$，
当 $\eta$ 非分歧且形如 $\eta(a)=|a|^{in\sigma}$ 时，至多在 $s=-i\sigma,\,1-i\sigma$ 处有单极点。
作为 $g$ 的函数，它光滑且具中等增长；作为 $s$ 的函数，它在竖直带中有界
（在可能极点之外），且对紧集中的 $g$ 一致。此外有函数方程
\[
E(g,\Phi;s,\eta)=E(g^\iota,\widehat{\Phi};1-s,\eta^{-1}),
\]
其中 $g^\iota={}^tg^{-1}$。
注意在中心之下，Eisenstein 级数按中心特征 $\eta^{-1}$ 变换。

现在回到 Euler 积分。设 $\pi$ 与 $\pi'$ 是不可约尖点表示，
其中心特征为 $\omega$ 与 $\omega'$。置 $\eta=\omega\omega'$。
则对每对尖点形式 $\varphi\in V_\pi$、$\varphi'\in V_{\pi'}$ 及每个
Schwartz--Bruhat 函数 $\Phi\in\mathcal{S}(\A^n)$，置
\[
I(s;\varphi,\varphi',\Phi)=\int_{Z_n(\A)\GL_n(k)\backslash\GL_n(\A)}
\varphi(g)\varphi'(g)E(g,\Phi;s,\eta)\,dg.
\]
由于两个尖点形式在 $Z_n(\A)\GL_n(k)\backslash\GL_n(\A)$ 上速降，
而 Eisenstein 级数只有中等增长，该积分对一切不在 Eisenstein 级数极点处的 $s$
绝对收敛，从而是亚纯的。它在极点之外的竖直带中有界，并满足函数方程
\[
I(s;\varphi,\varphi',\Phi)=I(1-s;\widetilde{\varphi},\widetilde{\varphi}',\widehat{\Phi}),
\]
此即 Eisenstein 级数函数方程的推论；此处仍有
$\widetilde{\varphi}(g)=\varphi(g^\iota)=\varphi(w_n{}^tg^{-1})\in V_{\widetilde{\pi}}$，
$\widetilde{\varphi}'$ 同理。

除非 $\eta(a)=\omega(a)\omega'(a)=|a|^{in\sigma}$ 非分歧，这些积分都是整的。
在那种情形，$s=-i\sigma$ 处的留数为
\[
\operatorname*{Res}_{s=-i\sigma}I(s;\varphi,\varphi',\Phi)
=-c\Phi(0)\int_{Z_n(\A)\GL_n(k)\backslash\GL_n(\A)}
\varphi(g)\varphi'(g)|\det(g)|^{-i\sigma}\,dg,
\]
而在 $s=1-i\sigma$ 处可把留数写成
\[
\operatorname*{Res}_{s=1-i\sigma}I(s;\varphi,\varphi',\Phi)
=c\widehat{\Phi}(0)\int_{Z_n(\A)\GL_n(k)\backslash\GL_n(\A)}
\widetilde{\varphi}(g)\widetilde{\varphi}'(g)|\det(g)|^{i\sigma}\,dg.
\]
因此这些留数定义了 $\pi$ 与 $\pi'\otimes|\det|^{-i\sigma}$ 之间的 $\GL_n(\A)$-不变配对，
等价地，也定义了 $\widetilde{\pi}$ 与 $\widetilde{\pi}'\otimes|\det|^{i\sigma}$ 之间的配对。
于是留数非零的必要条件是 $\pi\cong\widetilde{\pi}'\otimes|\det|^{i\sigma}$；
而在这种情形，确实可以选取 $\varphi$、$\varphi'$、$\Phi$ 使留数不为零。

还须验证积分是 Euler 乘积。为此取该积分，用定义替换 Eisenstein 级数并展开，得
\[
I(s;\varphi,\varphi',\Phi)=\int_{P_n(k)\backslash\GL_n(\A)}
\varphi(g)\varphi'(g)\Phi(e_ng)|\det(g)|^s\,dg.
\]
再用 Fourier 展开替换 $\varphi$，如前展开，得
\[
\begin{aligned}
I(s;\varphi,\varphi',\Phi)&=\Psi(s;W_\varphi,W'_{\varphi'},\Phi)\\
&=\int_{N_n(\A)\backslash\GL_n(\A)}
W_\varphi(g)W'_{\varphi'}(g)\Phi(e_ng)|\det(g)|^s\,dg.
\end{aligned}
\]
该式对 $\operatorname{Re}(s)\gg0$ 收敛。

往下，假定 $\varphi$、$\varphi'$、$\Phi$ 在同构
$\pi\cong\otimes'\pi_v$、$\pi'\cong\otimes'\pi'_v$、$\mathcal{S}(\A^n)\cong\otimes'\mathcal{S}(k_v^n)$
下都是可分解张量，于是
$W_\varphi(g)=\prod_vW_{\varphi_v}(g_v)$，
$W'_{\varphi'}(g)=\prod_vW'_{\varphi'_v}(g_v)$，
$\Phi(g)=\prod_v\Phi_v(g_v)$。
又因积分区域也自然分解，可把最后的积分拆成 Euler 乘积，写作
\[
\Psi(s;W_\varphi,W'_{\varphi'},\Phi)
=\prod_v\Psi_v(s;W_{\varphi_v},W'_{\varphi'_v},\Phi_v),
\]
其中
\[
\Psi_v(s;W_{\varphi_v},W'_{\varphi'_v},\Phi_v)
=\int_{N_n(k_v)\backslash\GL_n(k_v)}
W_{\varphi_v}(g_v)W'_{\varphi'_v}(g_v)\Phi_v(e_ng_v)|\det(g_v)|^s\,dg_v,
\]
仍由局部规范估计对 $\operatorname{Re}(s)\gg0$ 收敛。
我们再一次看到，Euler 分解是 Whittaker 模型唯一性的直接推论。

\textbf{定理 2.2.}\quad 设 $\varphi\in V_\pi$、$\varphi'\in V_{\pi'}$ 是 $\GL_n(\A)$
上的尖点形式，$\Phi\in\mathcal{S}(\A^n)$。则积分族 $I(s;\varphi,\varphi',\Phi)$
定义 $s$ 的亚纯函数，在极点之外的竖直带中有界。
可能的极点全为单极点，且仅当 $\pi\cong\widetilde{\pi}'\otimes|\det|^{i\sigma}$
（$\sigma$ 实）时出现，位置为 $s=-i\sigma$ 与 $s=1-i\sigma$，留数如上。
它们满足函数方程
\[
I(s;\varphi,\varphi',\Phi)=I(1-s;\widetilde{W}_\varphi,\widetilde{W}'_{\varphi'},\widehat{\Phi}).
\]
此外，若 $\varphi$、$\varphi'$、$\Phi$ 可分解，则积分是 Euler 乘积，且有
\[
I(s;\varphi,\varphi',\Phi)=\prod_v\Psi_v(s;W_{\varphi_v},W'_{\varphi'_v},\Phi_v),
\]
收敛绝对且对 $\operatorname{Re}(s)\gg0$ 一致。

顺带说明，函数方程的右端也可展开为
\[
\begin{aligned}
I(1-s;\widetilde{\varphi},\widetilde{\varphi}',\widehat{\Phi})
&=\int_{N_n(\A)\backslash\GL_n(\A)}
\widetilde{W}_\varphi(g)\widetilde{W}'_{\varphi'}(g)\widehat{\Phi}(e_ng)|\det(g)|^{1-s}\,dg\\
&=\prod_v\Psi_v(1-s;\widetilde{W}_{\varphi_v},\widetilde{W}'_{\varphi'_v},\widehat{\Phi}_v),
\end{aligned}
\]
对 $\operatorname{Re}(s)\ll0$ 收敛。
\section{局部 \texorpdfstring{$L$}{L} 函数}

若 $(\pi,V_\pi)$ 是 $\GL_n(\A)$ 的尖点表示，$(\pi',V_{\pi'})$ 是 $\GL_m(\A)$ 的尖点表示，
则已给数对 $(\pi,\pi')$ 配上一族 Euler 积分 $\{I(s;\varphi,\varphi')\}$
（当 $m=n$ 时为 $\{I(s;\varphi,\varphi',\Phi)\}$），
并经 Euler 分解给每个位置 $v$ 配上一族局部积分 $\{\Psi_v(s;W_v,W'_v)\}$
（或 $\{\Psi_v(s;W_v,W'_v,\Phi_v)\}$），它们附着于局部分量对 $(\pi_v,\pi'_v)$。
本节要通过这族局部积分，给这样一对局部表示配上局部 $L$ 函数（或局部 Euler 因子）
$L(s,\pi_v\times\pi'_v)$，并分析其基本性质，包括局部函数方程。

\subsection{非阿基米德局部因子}

本节中 $k$ 表非阿基米德局部域。以 $\mathfrak{o}$ 记 $k$ 的整数环，
$\mathfrak{p}$ 记 $\mathfrak{o}$ 的唯一素理想，并固定 $\mathfrak{p}$ 的一个生成元 $\varpi$。
以 $q$ 记 $k$ 的剩余域次数，故 $q=|\mathfrak{o}/\mathfrak{p}|=|\varpi|^{-1}$。
固定 $k$ 的一个非平凡连续加法特征 $\psi$。
此时 $(\pi,V_\pi)$ 与 $(\pi',V_{\pi'})$ 分别是 $\GL_n(k)$ 与 $\GL_m(k)$ 的
不可约容许酉一般表示中的光滑向量，尖点表示的局部分量正是如此。
以 $\omega$ 与 $\omega'$ 记其中心特征。本节的基本参考文献是 Jacquet、Piatetski-Shapiro
与 Shalika 的论文 [23]。

对每一对 Whittaker 函数 $W\in\mathcal{W}(\pi,\psi)$ 与 $W'\in\mathcal{W}(\pi',\psi^{-1})$，
在 $m=n$ 时还对每个 Schwartz--Bruhat 函数 $\Phi\in\mathcal{S}(k^n)$，
前面已定义了局部积分 $\Psi(s;W,W')$、$\widetilde{\Psi}(s;W,W')$ 与
$\Psi(s;W,W',\Phi)$，它们都对 $\operatorname{Re}(s)\gg0$ 收敛。
为下文记号方便，当 $m<n$ 时，对 $0\le j\le n-m-1$ 置
\[
\begin{aligned}
\Psi_j(s;W,W')={}&\int_{N_m(k)\backslash\GL_m(k)}\int_{M_{j,m}(k)}
W\left(\begin{pmatrix}h\\x&I_j\\&&I_{n-m-j}\end{pmatrix}\right)dx\\
&\cdot W'(h)|\det(h)|^{s-(n-m)/2}\,dh,
\end{aligned}
\]
于是 $\Psi(s;W,W')=\Psi_0(s;W,W')$，而
$\widetilde{\Psi}(s;W,W')=\Psi_{n-m-1}(s;W,W')$，两者仍对
$\operatorname{Re}(s)\gg0$ 绝对收敛。

我们需要弄清这些局部积分作为 $s$ 的函数属于何种类型。为此须先了解局部
Whittaker 函数。设 $W\in\mathcal{W}(\pi,\psi)$。由于 $W$ 光滑，存在紧开阔子群 $K$，
使对一切 $k\in K$ 有 $W(gk)=W(g)$。$W$ 又在左侧按 $\psi$ 作 $N_n(k)$-变换。
故由 $\GL_n(k)$ 的 Iwasawa 分解，只需了解环面上的广义 Whittaker 函数。
设 $\alpha_i$（$i=1,\ldots,n-1$）为 $\GL_n$ 的标准单根，即对
$a=\diag(a_1,\ldots,a_n)\in A_n(k)$ 有 $\alpha_i(a)=a_i/a_{i+1}$。
关于 Whittaker 函数渐近性的基本结果 [22] 是：存在只依赖 $\pi$ 的有限函数集
$X(\pi)=\{\chi_i\}$（定义在 $A_n(k)$ 上），使对每个 $W\in\mathcal{W}(\pi,\psi)$
都存在 Schwartz--Bruhat 函数 $\phi_i\in\mathcal{S}(k^{n-1})$，使得对一切满足
$a_n=1$ 的 $a\in A_n(k)$ 有
\[
W(a)=\sum_{X(\pi)}\chi_i(a)\phi_i(\alpha_1(a),\ldots,\alpha_{n-1}(a)).
\]
所谓 $A_n(k)$ 上的有限函数，指平移张成有限维向量空间的连续函数
[21, 22, 第 2.2 节]。对域 $k^\times$ 本身，有限函数由特征与赋值映射之幂的乘积张成。
这些出现在 Whittaker 函数零点附近渐近式中的有限函数集 $X(\pi)$，
来自分析 Jacquet 模 $\mathcal{W}(\pi,\psi)/\{\pi(n)W-W\mid n\in N_n\}$，
后者自然是一个 $A_n(k)$-模。
注意由于 Schwartz--Bruhat 函数的存在，只要某个单根 $\alpha_i(a)$ 变大，
Whittaker 函数便消失。

把 $W$ 与 $W'$ 的这些公式代入局部积分 $\Psi_j(s;W,W')$ 或 $\Psi(s;W,W',\Phi)$，
可得若干良好的推论 [22, 23]。
\begin{enumerate}
\item 每个积分对 $\operatorname{Re}(s)\gg0$ 收敛。当 $\pi$ 与 $\pi'$ 酉时（如我们所设），
它们在 $\operatorname{Re}(s)\ge1$ 时绝对收敛；当 $\pi$ 与 $\pi'$ 是缓增的（tempered）时，
在 $\operatorname{Re}(s)>0$ 时绝对收敛。
\item 每个积分定义 $q^{-s}$ 的有理函数，从而亚纯延拓到整个 $\mathbb{C}$。
\item 每个这样的有理函数可写成公分母形式，其分母只依赖有限函数集 $X(\pi)$ 与
$X(\pi')$，从而只依赖 $\pi$ 与 $\pi'$。
\end{enumerate}
\begin{annote}
这三条里真正起作用的是第三条：它说明所有局部积分的分母被同一个多项式控制，
因此它们落在 $\mathbb{C}(q^{-s})$ 的一个「分母有界」的子空间里。
术语上，$A_n(k)$ 上的有限函数就是 Jacquet 模的权，$X(\pi)$ 是这些权之集；
渐近展开写成 $\chi_i(a)\phi_i(\alpha_1(a),\ldots)$ 之后，$\phi_i$ 只依赖根坐标，
于是对 $a$ 的依赖被完全分离，积分才有机会化成 $q^{-s}$ 的有理函数。
\par\vspace{0.15em}{\raggedleft 【来源：Jacquet--Piatetski-Shapiro--Shalika, Amer.\ J.\ Math.\ 105 (1983), §2 与 §9（A）】\par}
\end{annote}

记 $\mathcal{I}_j(\pi,\pi')$ 为 $m<n$ 时局部积分 $\Psi_j(s;W,W')$ 的复线性张成，
记 $\mathcal{I}(\pi,\pi')$ 为 $m=n$ 时 $\Psi(s;W,W',\Phi)$ 的张成。
当 $m<n$ 时可以证明，用相当初等却不无繁琐的积分变形，所有理想
$\mathcal{I}_j(\pi,\pi')$ 都相同 [23]，故统一记此理想为 $\mathcal{I}(\pi,\pi')$。
这些子空间是 $\mathbb{C}(q^{-s})$ 中按 (3) 的意义「分母有界」的子空间。
事实上它们还有更多结构：每个 $\mathcal{I}(\pi,\pi')$ 都是 $\mathbb{C}(q^{-s})$ 中
分数 $\mathbb{C}[q^s,q^{-s}]$-理想。
由于 $\mathbb{C}[q^s,q^{-s}]$ 是主理想整环，每个分数理想 $\mathcal{I}(\pi,\pi')$
有单个生成元；又因这些分数理想都含 $1$，总可把生成元归一化为
$P_{\pi,\pi'}(q^{-s})^{-1}$ 的形状，其中多项式 $P(X)$ 满足 $P(0)=1$。
由此得到局部 $L$ 函数的定义：
\[
L(s,\pi\times\pi')=P_{\pi,\pi'}(q^{-s})^{-1}
\]
是局部积分族所成分数理想 $\mathcal{I}(\pi,\pi')$ 的归一化生成元。
若 $\pi'=\mathbf{1}$ 是 $\GL_1(k)$ 的平凡表示，则记 $L(s,\pi)=L(s,\pi\times\mathbf{1})$。

容易证明，理想 $\mathcal{I}(\pi,\pi')$ 与定义 Whittaker 模型所用的特征 $\psi$ 无关，
故 $L(s,\pi\times\pi')$ 与 $\psi$ 的选取无关，因而不出现在记号中。
又注意若 $\pi'=\chi$ 是 $\GL_1(\A)$ 的自守表示（特征），则有恒等式
$L(s,\pi\times\chi)=L(s,\pi\otimes\chi)$，其中 $\pi\otimes\chi$ 是 $\GL_n(\A)$ 在
$V_\pi$ 上由 $\pi\otimes\chi(g)\xi=\chi(\det(g))\pi(g)\xi$ 给出的表示。
我们在下述定理中总结以上内容。

\textbf{定理 3.1.}\quad 设 $\pi$ 与 $\pi'$ 如上。局部积分族构成
$\mathbb{C}(q^{-s})$ 中的分数 $\mathbb{C}[q^s,q^{-s}]$-理想 $\mathcal{I}(\pi,\pi')$，
其生成元即局部 $L$ 函数 $L(s,\pi\times\pi')$。

局部 $L$ 函数还有另一种有用的看法。
$L(s,\pi\times\pi')$ 是形如 $P(q^{-s})^{-1}$ 的函数中次数最小者，
其中 $P(X)$ 是满足 $P(0)=1$ 的多项式，且使对一切 $W\in\mathcal{W}(\pi,\psi)$、
$W'\in\mathcal{W}(\pi',\psi^{-1})$（必要时还有 $\Phi\in\mathcal{S}(k^n)$）商
\[
\frac{\Psi(s;W,W')}{L(s,\pi\times\pi')}
\quad\text{（相应地}\quad
\frac{\Psi(s;W,W',\Phi)}{L(s,\pi\times\pi')}\text{）}
\]
都是整函数。也就是说，$L(s,\pi\times\pi')$ 是由 $\mathcal{I}(\pi,\pi')$ 中函数的
极点所确定的标准 Euler 因子。
\begin{annote}
用分数理想来定义 $L$ 因子，是这一理论中最容易被误解的一步。
先验地，我们并不知道局部 $L$ 因子该长什么样；能确认的是局部积分落在某个
一维生成元所刻画的理想里，而生成元在归一化 $P(0)=1$ 之后唯一。
如此定义的 $L$ 因子自动满足「最小性」，即不能再去掉任何因子，
这正是函数方程中 $\varepsilon$ 因子能成为单项式的原因。
反过来，若想验证某处算出的 $L$ 因子正确，只需验证同一族积分给出的理想与之一致。
\par\vspace{0.15em}{\raggedleft 【来源：Jacquet--Piatetski-Shapiro--Shalika, Amer.\ J.\ Math.\ 105 (1983), §3；Cogdell, ICTP Lecture Notes, §3（A）】\par}
\end{annote}

应当注意，既然 $L$ 因子是理想 $\mathcal{I}(\pi,\pi')$ 的生成元，它特别属于
$\mathcal{I}(\pi,\pi')$。又因该理想由局部积分张成，总可写
\[
L(s,\pi\times\pi')=\sum_i\Psi(s;W_i,W'_i)
\quad\text{或}\quad
L(s,\pi\times\pi')=\sum_i\Psi(s;W_i,W'_i,\Phi_i),
\]
其中 $W_i\in\mathcal{W}(\pi,\psi)$、$W'_i\in\mathcal{W}(\pi',\psi^{-1})$ 取有限多个，
必要时还有 $\Phi_i\in\mathcal{S}(k^n)$。这一点在整体理论中要用到。

无论与 Tate 论文作类比，还是看相应的整体陈述，我们都应期待局部积分满足
局部函数方程。由整体积分的函数方程，应当期待当 $m<n$ 时它把
$\Psi(s;W,W')$ 与 $\widetilde{\Psi}(1-s;\rho(w_{n,m})\widetilde{W},\widetilde{W}')$
联系起来，而当 $m=n$ 时把 $\Psi(s;W,W',\Phi)$ 与
$\Psi(1-s;\widetilde{W},\widetilde{W}',\widehat{\Phi})$ 联系起来。
事实确是如此。这些函数方程来自把局部积分解释为一族以 $s$ 为参数的双线性型，
它们定义在 $\mathcal{W}(\pi,\psi)\times\mathcal{W}(\pi',\psi^{-1})$ 上；
当 $m=n$ 时则是 $\mathcal{W}(\pi,\psi)\times\mathcal{W}(\pi',\psi^{-1})
\times\mathcal{S}(k^n)$ 上的三线性型。
利用 Bruhat 理论可以分析这样的泛函，并证明除 $q^{-s}$ 的有限多个例外值外，
这样的双线性型是唯一的 [23]。于是得到下述局部函数方程。

\textbf{定理 3.2.}\quad 存在有理函数
$\gamma(s,\pi\times\pi',\psi)\in\mathbb{C}(q^{-s})$，使得当 $m<n$ 时
\[
\widetilde{\Psi}(1-s;\rho(w_{n,m})\widetilde{W},\widetilde{W}')
=\omega'(-1)^{n-1}\gamma(s,\pi\times\pi',\psi)\Psi(s;W,W'),
\]
而当 $m=n$ 时
\[
\Psi(1-s;\widetilde{W},\widetilde{W}',\widehat{\Phi})
=\omega'(-1)^{n-1}\gamma(s,\pi\times\pi',\psi)\Psi(s;W,W',\Phi),
\]
对一切 $W\in\mathcal{W}(\pi,\psi)$、$W'\in\mathcal{W}(\pi',\psi^{-1})$，
必要时对一切 $\Phi\in\mathcal{S}(k^n)$ 成立。

另一同等重要的局部因子，出现在局部 Langlands 对应的现今表述中 [2, 16, 19]，
即局部 $\varepsilon$ 因子，定义为比值
\[
\varepsilon(s,\pi\times\pi',\psi)
=\frac{\gamma(s,\pi\times\pi',\psi)L(s,\pi\times\pi')}
{L(1-s,\widetilde{\pi}\times\widetilde{\pi}')}.
\]
有了局部 $\varepsilon$ 因子，局部函数方程可写成下述形式：当 $m<n$ 时
\[
\frac{\widetilde{\Psi}(1-s;\rho(w_{n,m})\widetilde{W},\widetilde{W}')}
{L(1-s,\widetilde{\pi}\times\widetilde{\pi}')}
=\omega'(-1)^{n-1}\varepsilon(s,\pi\times\pi',\psi)
\frac{\Psi(s;W,W')}{L(s,\pi\times\pi')},
\]
而当 $m=n$ 时
\[
\frac{\Psi(1-s;\widetilde{W},\widetilde{W}',\widehat{\Phi})}
{L(1-s,\widetilde{\pi}\times\widetilde{\pi}')}
=\omega'(-1)^{n-1}\varepsilon(s,\pi\times\pi',\psi)
\frac{\Psi(s;W,W',\Phi)}{L(s,\pi\times\pi')}.
\]

按这种形式分析局部函数方程，可以证明 $\varepsilon(s,\pi\times\pi',\psi)$
是形如 $cq^{-fs}$ 的单项式。若只考虑单个 $\pi$，取 $\psi$ 无分歧，并记
$\varepsilon(s,\pi,\psi)=\varepsilon(0,\pi,\psi)q^{-f(\pi)s}$，
则 [24] 证明 $f(\pi)$ 是非负整数，$f(\pi)=0$ 当且仅当 $\pi$ 无分歧；
一般地，整数 $f(\pi)$ 就是 $\pi$ 的\emph{导子}（conductor）。

下面转向局部 $L$ 函数的计算。第一个要考虑的情形是 $\pi$ 与 $\pi'$ 都无分歧。
由于它们假定是一般的，二者都是 Borel 子群的无分歧特征的全诱导表示 [49]。
故可写 $\pi\cong\operatorname{Ind}_{B_n}^{\GL_n}(\mu_1\otimes\cdots\otimes\mu_n)$ 与
$\pi'\cong\operatorname{Ind}_{B_m}^{\GL_m}(\mu'_1\otimes\cdots\otimes\mu'_m)$，
其中 $\mu_i$ 与 $\mu'_j$ 是 $k^\times$ 的无分歧特征。
无分歧表示的 Satake 参数化给每个这样的表示配上半单共轭类
$[A_\pi]\in\GL_n(\mathbb{C})$ 与 $[A_{\pi'}]\in\GL_m(\mathbb{C})$，即
\[
A_\pi=\begin{pmatrix}\mu_1(\varpi)&&\\&\ddots&\\&&\mu_n(\varpi)\end{pmatrix},
\qquad
A_{\pi'}=\begin{pmatrix}\mu'_1(\varpi)&&\\&\ddots&\\&&\mu'_m(\varpi)\end{pmatrix}.
\]
（回忆 $\varpi$ 是 $k$ 的一致化子，即 $\mathfrak{p}$ 的生成元。）
在 Whittaker 模型中有唯一的归一化 $K=\GL(\mathfrak{o})$-不变 Whittaker 函数
$W_{\mathrm{o}}\in\mathcal{W}(\pi,\psi)$ 与
$W'_{\mathrm{o}}\in\mathcal{W}(\pi',\psi^{-1})$，归一化为
$W_{\mathrm{o}}(e)=W'_{\mathrm{o}}(e)=1$。
Shintani [43] 给出了 $W_{\mathrm{o}}$ 用 Satake 参数 $A_\pi$ 表示的显式公式。
利用该公式，便得到此情形局部 $L$ 因子的如下显式计算。

\textbf{定理 3.3.}\quad 若 $\pi$、$\pi'$ 与 $\psi$ 都无分歧，则
\[
L(s,\pi\times\pi')=\det(I-q^{-s}A_\pi\otimes A_{\pi'})^{-1}
=\begin{cases}
\Psi(s;W_{\mathrm{o}},W'_{\mathrm{o}}),&m<n,\\
\Psi(s;W_{\mathrm{o}},W'_{\mathrm{o}},\Phi_{\mathrm{o}}),&m=n,
\end{cases}
\]
且 $\varepsilon(s,\pi\times\pi',\psi)\equiv1$。
\begin{annote}
定理 3.3 值得与经典情形对照：当 $n=m=1$ 时，
$\det(I-q^{-s}\mu(\varpi)\mu'(\varpi))^{-1}$ 就是通常的
$(1-\mu\mu'(\varpi)q^{-s})^{-1}$，即两个 Hecke 特征之积的局部 zeta 因子。
一般 $n,m$ 时，$\det(I-q^{-s}A_\pi\otimes A_{\pi'})^{-1}$
是 $nm$ 个几何级数之积，展开就是 Rankin--Selberg 局部因子的标准形状。
定理中把 $L$ 因子与单个积分值等同起来，这一点只对无分歧表示成立；
有分歧时积分不再等于 $L$，而只落在以 $L$ 为生成元的理想里。
\par\vspace{0.15em}{\raggedleft 【来源：Shintani, \textit{On an explicit formula for class-1 ``Whittaker functions'' on $\mathrm{GL}_n$ over $\mathfrak{p}$-adic fields}, Proc.\ Japan Acad.\ 52 (1976), 180--182；Casselman--Shalika, Compositio Math.\ 41 (1980)（A）】\par}
\end{annote}

另一基本情形是 $\pi$ 与 $\pi'$ 都是超尖点的。此计算须用表示的 Kirillov 模型
来分析局部积分 [11, 8]。

\textbf{定理 3.4.}\quad 若 $\pi$ 与 $\pi'$ 都是（酉）超尖点的，则当 $m<n$ 时
$L(s,\pi\times\pi')\equiv1$；而当 $m=n$ 时有
\[
L(s,\pi\times\pi')=\prod(1-\alpha q^{-s})^{-1},
\]
乘积取遍所有满足 $\widetilde{\pi}\cong\pi'\otimes|\det|^{s_0}$ 的 $\alpha=q^{s_0}$。

其余情形只能借助 Bernstein--Zelevinsky 对 $\GL_n(k)$ 一般表示的分类 [49]。
一切一般表示都可实现为由超尖点表示抛物诱导所得表示的特定合成因子。
可以把诱导表示的 Whittaker 函数用诱导数据的 Whittaker 函数表出（如 [23] 所做），
也可以按 [8] 用 $\pi$ 与 $\pi'$ 的导数的拟不变配对来分析局部积分的极点，
从而把 $L(s,\pi\times\pi')$ 用超尖点表示对的 $L$ 函数算出。显式公式请参看
上述论文或 [28]。

作为本节的结尾，我们提两条关于 $\gamma$ 因子的结果。其一用于一般情形下
$L$ 因子的计算，即 $\gamma$ 因子的\emph{乘性} [23]。
其二即 $\gamma$ 因子的\emph{稳定性} [26]。这两条在把逆定理用于提升时都必需。

\textbf{$\gamma$ 因子的乘性.}\quad 若 $\pi=\operatorname{Ind}(\pi_1\otimes\pi_2)$，
其中 $\pi_i$ 是 $\GL_{r_i}(k)$ 的不可约容许表示，则
\[
\gamma(s,\pi\times\pi',\psi)=\gamma(s,\pi_1\times\pi',\psi)\gamma(s,\pi_2\times\pi',\psi),
\]
对 $\pi'$ 亦然。此外 $L(s,\pi\times\pi')^{-1}$ 整除
$[L(s,\pi_1\times\pi')L(s,\pi_2\times\pi')]^{-1}$。

\textbf{$\gamma$ 因子的稳定性.}\quad 若 $\pi_1$ 与 $\pi_2$ 是 $\GL_n(k)$ 的两个
不可约容许一般表示，且有相同的中心特征，则对每个充分分歧的 $\GL_1(k)$ 特征 $\eta$ 有
\[
\gamma(s,\pi_1\times\eta,\psi)=\gamma(s,\pi_2\times\eta,\psi),
\qquad
L(s,\pi_1\times\eta)=L(s,\pi_2\times\eta)\equiv1.
\]
更一般地，若再有 $\pi'$ 是 $\GL_m(k)$ 的不可约一般表示，则对一切充分分歧的
$\GL_1(k)$ 特征 $\eta$ 有
\[
\gamma(s,(\pi_1\otimes\eta)\times\pi',\psi)
=\gamma(s,(\pi_2\otimes\eta)\times\pi',\psi),
\qquad
L(s,(\pi_1\otimes\eta)\times\pi')=L(s,(\pi_2\otimes\eta)\times\pi')\equiv1.
\]

\subsection{阿基米德局部因子}

阿基米德局部因子的处理与非阿基米德情形在许多方面平行，但有几处显著差别。
这方面的主要工作是 Jacquet 与 Shalika 的 [27]，我们大体依循之而不再逐一标注，
此外还用到 [25] 中的阿基米德部分。

阿基米德理论发展中的一处显著差别是：理论展开之时，局部 Langlands 对应已经就位 [30]。
用通常的 Langlands 分类，这一对应可以写得很显式。
于是 $\pi$ 对应 $k$ 的 Weil 群 $W_k$ 的一个 $n$ 维半单表示
$\rho=\rho_\pi$，$\pi'$ 对应 $W_k$ 的一个 $m$ 维半单表示 $\rho'=\rho'_{\pi'}$。
则 $\rho\otimes\rho'$ 是 $W_k$ 的 $nm$ 维表示，而此 Weil 群表示配有
Artin--Weil 的 $L$ 因子与 $\varepsilon$ 因子 [47]，记作 $L(s,\rho\otimes\rho')$
与 $\varepsilon(s,\rho\otimes\rho',\psi)$。Jacquet 与 Shalika 实质上\emph{定义}
\[
L(s,\pi\times\pi')=L(s,\rho_\pi\otimes\rho'_{\pi'}),
\qquad
\varepsilon(s,\pi\times\pi',\psi)=\varepsilon(s,\rho_\pi\otimes\rho'_{\pi'},\psi),
\]
然后置
\[
\gamma(s,\pi\times\pi',\psi)
=\frac{\varepsilon(s,\pi\times\pi',\psi)L(1-s,\widetilde{\pi}\times\widetilde{\pi}')}
{L(s,\pi\times\pi')}.
\]
\begin{annote}
阿基米德情形与非阿基米德情形最大的不同，就在 $L$ 因子的来源：
非阿基米德时 $L$ 由局部积分的理想反推出来，阿基米德时则直接搬用
Artin--Weil 因子，再回头验证局部积分确实给出同一个理想。
这也解释了为什么 $\Gamma$ 因子在阿基米德情形是「标准型」的乘积，
而非阿基米德情形只是一串 $q^{-s}$ 的多项式。
$\rho_\pi$ 由 Langlands 分类给出，半单性来自 $\pi$ 的完全可约性；
若 $\pi$ 缓增，则 $\rho_\pi$ 还是酉的。
\par\vspace{0.15em}{\raggedleft 【来源：Jacquet--Shalika, \textit{Rankin--Selberg convolutions: Archimedean theory}, 1990, 125--207；Langlands, \textit{On the classification of irreducible representations of real algebraic groups}, 1989（A）】\par}
\end{annote}

他们进而证明这些函数与局部积分有预期的关系。为此定义 $\mathcal{M}(\pi\times\pi')$
为一切亚纯函数 $\phi(s)$ 所成之空间，其性质为：若 $P(s)$ 是使
$P(s)L(s,\pi\times\pi')$ 在竖直带 $S[a,b]=\{s\mid a\le\operatorname{Re}(s)\le b\}$
上全纯的多项式，则 $P(s)\phi(s)$ 在 $S[a,b]$ 上有界。
特别地，若 $\phi\in\mathcal{M}(\pi\times\pi')$，则商
$\phi(s)L(s,\pi\times\pi')^{-1}$ 是整函数。
然后按非阿基米德情形的方式定义局部积分 $\Psi_j(s;W,W')$ 与
$\Psi(s;W,W',\Phi)$，其中 $W\in\mathcal{W}(\pi,\psi)$、
$W'\in\mathcal{W}(\pi',\psi^{-1})$、$\Phi\in\mathcal{S}(k^n)$，
并用 [23] 中非阿基米德情形的直接类比方法来分析它们。

\textbf{定理 3.5.}\quad 积分 $\Psi_j(s;W,W')$ 或 $\Psi(s;W,W',\Phi)$ 延拓为
$s$ 的亚纯函数，且落在 $\mathcal{M}(\pi\times\pi')$ 中。
特别地，比值
\[
e_j(s;W,W')=\frac{\Psi_j(s;W,W')}{L(s,\pi\times\pi')},
\qquad
e(s;W,W',\Phi)=\frac{\Psi(s;W,W',\Phi)}{L(s,\pi\times\pi')}
\]
是整函数，且在竖直带中有界。

此陈述比非阿基米德情形「延拓」与「分母有界」两条包含更多内容。
既然它把「分母」指定为 $L$ 因子 $L(s,\pi\times\pi')^{-1}$，
就与实际计算局部积分的极点绑在一起。
事实上，Jacquet 与 Shalika [27] 相当大的一部分篇幅，正用于同时证明此事
与局部函数方程。

\textbf{定理 3.6.}\quad 有局部函数方程
\[
\Psi_{n-m-j-1}(1-s;\rho(w_{n,m})\widetilde{W},\widetilde{W}')
=\omega'(-1)^{n-1}\gamma(s,\pi\times\pi',\psi)\Psi_j(s;W,W'),
\]
或
\[
\Psi(1-s;\widetilde{W},\widetilde{W}',\widehat{\Phi})
=\omega'(-1)^{n-1}\gamma(s,\pi\times\pi',\psi)\Psi(s;W,W',\Phi).
\]

此时所缺的是 $L$ 因子的「最小性」陈述。
也就是说，我们已知 $L(s,\pi\times\pi')$ 是标准的阿基米德 Euler 因子，
即标准型 $\Gamma$ 函数之积。我们还知道，所有局部积分的极点都含在
$L$ 因子的极点之中，且计及重数。但尚未证明 $L$ 因子不能有多余的极点。
尤其不知道能否把局部 $L$ 函数表成局部积分的有限线性组合。
为此 Jacquet 与 Shalika 扩大了允许的局部积分空间。
设 $\Lambda$ 与 $\Lambda'$ 分别是 $V_\pi$ 与 $V_{\pi'}$ 上与 Whittaker 模型
$\mathcal{W}(\pi,\psi)$、$\mathcal{W}(\pi',\psi^{-1})$ 相伴的 Whittaker 泛函。
则 $\widehat{\Lambda}=\Lambda\otimes\Lambda'$ 定义了代数张量积
$V_\pi\otimes V_{\pi'}$ 上的连续线性泛函，并连续延拓到拓扑张量积
$V_{\pi\widehat{\otimes}\pi'}=V_\pi\widehat{\otimes}V_{\pi'}$，
后者视为 $\GL_n(k)\times\GL_m(k)$ 的表示。现在置
\[
\mathcal{W}(\pi\widehat{\otimes}\pi',\psi)
=\left\{W(g,h)=\widehat{\Lambda}(\pi(g)\otimes\pi'(h)\xi)\ \middle|\
\xi\in V_{\pi\widehat{\otimes}\pi'}\right\}.
\]

于是 $\mathcal{W}(\pi\widehat{\otimes}\pi',\psi)$ 含代数张量积
$\mathcal{W}(\pi,\psi)\otimes\mathcal{W}(\pi',\psi^{-1})$，且等于拓扑张量积。
现在可把一切局部积分延拓到 $\mathcal{W}(\pi\widehat{\otimes}\pi',\psi)$，只须置
\[
\Psi_j(s;W)=\iint W\left(\begin{pmatrix}h\\x&I_j\\&&I_{n-m-j}\end{pmatrix},h\right)
dx\,|\det(h)|^{s-(n-m)/2}\,dh
\]
与
\[
\Psi(s;W,\Phi)=\int W(g,g)\Phi(e_ng)|\det(g)|^s\,dh,
\]
其中 $W\in\mathcal{W}(\pi\widehat{\otimes}\pi',\psi)$。
由于局部积分对拓扑张量积上的拓扑连续，以上一切结论仍然成立，
特别是收敛性陈述、局部函数方程，以及这些积分延拓为
$\mathcal{M}(\pi\times\pi')$ 中函数这一事实。
此时置 $\mathcal{I}_j(\pi,\pi')=\{\Psi_j(s;W)\mid
W\in\mathcal{W}(\pi\widehat{\otimes}\pi')\}$，并令 $\mathcal{I}(\pi,\pi')$ 为
$\{\Psi(s;W,\Phi)\mid W\in\mathcal{W}(\pi\widehat{\otimes}\pi'),\
\Phi\in\mathcal{S}(k^n)\}$ 的张成。
当 $m<n$ 时，空间 $\mathcal{I}_j(\pi,\pi')$ 与 $j$ 无关，故仍记之为
$\mathcal{I}(\pi,\pi')$。这些空间仍与 $\psi$ 无关。于是由上已知
$\mathcal{I}(\pi,\pi')\subset\mathcal{M}(\pi\times\pi')$。
[27] 余下的篇幅便用于证明 $\mathcal{I}(\pi,\pi')=\mathcal{M}(\pi\times\pi')$。

当 $m=n-1$ 或 $m=n$ 时，Stade [44, 45] 与 Jacquet、Shalika（见 [9]）证明，
确实可以把局部 $L$ 函数表成只涉及 $\mathcal{W}(\pi,\psi)$ 与
$\mathcal{W}(\pi',\psi^{-1})$ 中 $K$-有限函数的积分的有限线性组合，
也就是说，不必过渡到 $\mathcal{W}(\pi,\psi)\otimes\mathcal{W}(\pi',\psi^{-1})$
的完备化。

\section{整体 \texorpdfstring{$L$}{L} 函数}

仍设 $k$ 为整体域，$\A$ 为其阿代尔环，并固定 $\A$ 的一个在 $k$ 上平凡的非平凡
连续加法特征 $\psi=\otimes\psi_v$。

设 $(\pi,V_\pi)$ 是 $\GL_n(\A)$ 的尖点表示，$(\pi',V_{\pi'})$ 是 $\GL_m(\A)$
的尖点表示。由于二者不可约，它们有受限张量积分解 $\pi\cong\otimes'\pi_v$ 与
$\pi'\cong\otimes'\pi'_v$，其中 $(\pi_v,V_{\pi_v})$ 与 $(\pi'_v,V_{\pi'_v})$
分别是 $\GL_n(k_v)$ 与 $\GL_m(k_v)$ 的不可约容许光滑酉一般表示 [10, 13]。
设 $\omega=\otimes'\omega_v$ 与 $\omega'=\otimes'\omega'_v$ 为其中心特征，
二者都是 $k^\times\backslash\A^\times$ 的连续特征。
设 $S$ 是 $k$ 的有限位置集，含阿基米德位置集 $S_\infty$，
使得对一切 $v\notin S$，$\pi_v$、$\pi'_v$ 与 $\psi_v$ 都无分歧。

对 $k$ 的每个位置 $v$，已定义了局部因子 $L(s,\pi_v\times\pi'_v)$ 与
$\varepsilon(s,\pi_v\times\pi'_v,\psi_v)$。于是至少可以形式地定义
\[
L(s,\pi\times\pi')=\prod_vL(s,\pi_v\times\pi'_v),
\qquad
\varepsilon(s,\pi\times\pi')=\prod_v\varepsilon(s,\pi_v\times\pi'_v,\psi_v).
\]

定义 $L$ 函数的乘积对 $\operatorname{Re}(s)\gg0$ 绝对收敛。
对 $\varepsilon$ 因子，当 $v\notin S$ 时
$\varepsilon(s,\pi_v\times\pi'_v,\psi_v)\equiv1$，故该乘积实为有限乘积，
不存在收敛问题。$\varepsilon(s,\pi\times\pi')$ 与 $\psi$ 无关这一事实，
既可以按 [5, 引理 2.1] 分析局部 $\varepsilon$ 因子随 $\psi$ 的变化来验证，
也可由下文所述的整体函数方程推出。

\subsection{基本解析性质}

我们的首要目标是证明这些 $L$ 函数具有良好的解析性质。

\textbf{定理 4.1.}\quad 整体 $L$ 函数 $L(s,\pi\times\pi')$ 是「好的」，意即
\begin{enumerate}
\item $L(s,\pi\times\pi')$ 亚纯延拓到整个 $\mathbb{C}$；
\item 延拓后的函数在竖直带中有界（在极点之外）；
\item 满足函数方程
\[
L(s,\pi\times\pi')=\varepsilon(s,\pi\times\pi')
L(1-s,\widetilde{\pi}\times\widetilde{\pi}').
\]
\end{enumerate}

为此须把 $L$ 函数与整体积分联系起来。
先看延拓。当 $m<n$ 时，对每个 $\varphi\in V_\pi$ 与 $\varphi'\in V_{\pi'}$，
已知积分 $I(s;\varphi,\varphi')$ 对一切 $s$ 收敛。
由第 2 节的展开与第 3 节的局部计算，已知当 $\operatorname{Re}(s)\gg0$ 时，
对适当选取的 $\varphi$ 与 $\varphi'$ 有
\[
\begin{aligned}
I(s;\varphi,\varphi')&=\prod_v\Psi_v(s;W_{\varphi_v},W'_{\varphi'_v})
=\left(\prod_{v\in S}\Psi_v(s;W_{\varphi_v},W'_{\varphi'_v})\right)L^S(s,\pi\times\pi')\\
&=\left(\prod_{v\in S}\frac{\Psi_v(s;W_{\varphi_v},W'_{\varphi'_v})}
{L(s,\pi_v\times\pi'_v)}\right)L(s,\pi\times\pi')\\
&=\left(\prod_{v\in S}e_v(s;W_{\varphi_v},W'_{\varphi'_v})\right)L(s,\pi\times\pi').
\end{aligned}
\]
已知每个比值 $e_v(s;W_v,W'_v)$ 都是整函数，故 $L(s,\pi\times\pi')$ 有亚纯延拓。
若 $m=n$，则对适当的 $\varphi\in V_\pi$、$\varphi'\in V_{\pi'}$ 与
$\Phi\in\mathcal{S}(\A^n)$，同样有
\[
I(s;\varphi,\varphi',\Phi)=\left(\prod_{v\in S}e_v(s;W_{\varphi_v},W'_{\varphi'_v},\Phi_v)\right)
L(s,\pi\times\pi'),
\]
故 $L(s,\pi\times\pi')$ 也有亚纯延拓。

下面转向函数方程。它将由整体积分的函数方程与局部函数方程推出。
只考虑 $m<n$ 的情形，另一情形完全类似。
整体积分的函数方程就是
\[
I(s;\varphi,\varphi')=\widetilde{I}(1-s;\widetilde{\varphi},\widetilde{\varphi}').
\]
对适当的 $\varphi$ 与 $\varphi'$ 仍有
\[
I(s;\varphi,\varphi')=\left(\prod_{v\in S}
\frac{\Psi(s;W_v,W'_v)}{L(s,\pi_v\times\pi'_v)}\right)L(s,\pi\times\pi'),
\]
而在另一侧有
\[
\widetilde{I}(1-s;\widetilde{\varphi},\widetilde{\varphi}')
=\left(\prod_{v\in S}
\frac{\widetilde{\Psi}(1-s;\rho(w_{n,m})\widetilde{W}_v,\widetilde{W}'_v)}
{L(1-s,\widetilde{\pi}_v\times\widetilde{\pi}'_v)}\right)
L(1-s,\widetilde{\pi}\times\widetilde{\pi}').
\]
然而由局部函数方程，对每个 $v\in S$ 有
\[
\frac{\widetilde{\Psi}(1-s;\rho(w_{n,m})\widetilde{W}_v,\widetilde{W}'_v)}
{L(1-s,\widetilde{\pi}_v\times\widetilde{\pi}'_v)}
=\omega'_v(-1)^{n-1}\varepsilon(s,\pi_v\times\pi'_v,\psi_v)
\frac{\Psi(s;W_v,W'_v)}{L(s,\pi_v\times\pi'_v)}.
\]
合并以上各式，得
\[
L(s,\pi\times\pi')=\left(\prod_{v\in S}\omega'_v(-1)^{n-1}
\varepsilon(s,\pi_v\times\pi'_v,\psi_v)\right)
L(1-s,\widetilde{\pi}\times\widetilde{\pi}').
\]
现在对 $v\notin S$，已知 $\pi'_v$ 无分歧，故 $\omega'_v(-1)=1$，
且 $\varepsilon(s,\pi_v\times\pi'_v,\psi_v)\equiv1$。因此
\[
\prod_{v\in S}\omega'_v(-1)^{n-1}\varepsilon(s,\pi_v\times\pi'_v,\psi_v)
=\varepsilon(s,\pi\times\pi'),
\]
于是确实有
\[
L(s,\pi\times\pi')=\varepsilon(s,\pi\times\pi')
L(1-s,\widetilde{\pi}\times\widetilde{\pi}').
\]
注意这也说明 $\varepsilon(s,\pi\times\pi')$ 与 $\psi$ 无关。
\begin{annote}
这一推导把整体的困难全部推给了局部：整体函数方程之所以成立，
只因为每个 $v$ 处的局部函数方程乘起来之后，各个 $\gamma$ 因子与
$L$ 因子恰好相消，只剩下 $\varepsilon$。
因此第 3 节局部理论中「$L$ 是理想的生成元」这一句才是真正的关键；
若局部 $L$ 因子取得不准，乘积会多出无穷多个 $v\notin S$ 的因子，
整体函数方程便无从谈起。
另外注意 $\omega'_v(-1)$ 只对非阿基米德且分歧的 $v$ 才可能不为 $1$，
有限乘积中它被吸收进 $\varepsilon$ 的常数因子。
\par\vspace{0.15em}{\raggedleft 【来源：Jacquet--Shalika, Amer.\ J.\ Math.\ 103 (1981), 499--588；Cogdell, ICTP Lecture Notes, §4（A）】\par}
\end{annote}

下面转向竖直带中的有界性。对整体积分 $I(s;\varphi,\varphi')$ 或
$I(s;\varphi,\varphi',\Phi)$，有界性直接来自绝对收敛。
对 $L$ 函数本身，范式如下。对每个有限位置 $v\in S$，已知存在
$W_{v,i}$、$W'_{v,i}$，必要时还有 $\Phi_{v,i}$，使得
\[
L(s,\pi_v\times\pi'_v)=\sum_i\Psi(s;W_{v,i},W'_{v,i})
\quad\text{或}\quad
L(s,\pi_v\times\pi'_v)=\sum_i\Psi(s;W_{v,i},W'_{v,i},\Phi_{v,i}).
\]
若 $m=n-1$ 或 $m=n$，则由 Stade 以及 Jacquet 与 Shalika 的工作，
已知对 $v\in S_\infty$ 也有类似陈述。故当 $m=n-1$ 或 $m=n$ 时，
存在整体的 $\varphi_i$、$\varphi'_i$，必要时还有 $\Phi_i$，使得
\[
L(s,\pi\times\pi')=\sum_iI(s;\varphi_i,\varphi'_i)
\quad\text{或}\quad
L(s,\pi\times\pi')=\sum_iI(s;\varphi_i,\varphi'_i,\Phi_i).
\]
于是 $L$ 函数在竖直带中的有界性由整体积分的相应性质推出。

然而若 $m<n-1$，则在那些 $v\in S_\infty$ 处，我们只知道存在
$W_v\in\mathcal{W}(\pi_v\widehat{\otimes}\pi'_v,\psi_v)
=\mathcal{W}(\pi_v,\psi_v)\widehat{\otimes}\mathcal{W}(\pi'_v,\psi_v^{-1})$
或有限多个这样的函数 $W_{v,i}$ 与 $\Phi_{v,i}$，使得
\[
L(s,\pi_v\times\pi'_v)=I(s;W_v)
\quad\text{或}\quad
L(s,\pi_v\times\pi'_v)=\sum_iI(s;W_{v,i},\Phi_{v,i}).
\]
要使上述范式对 $m<n-1$ 也奏效，就需要重做 $V_\pi\otimes V_{\pi'}$ 中尖点形式的
整体 Euler 积分理论。后者自然是 $\GL_n(\A)\times\GL_m(\A)$ 的某个不可约酉尖点
表示的光滑向量空间。于是需要把整体积分理论平行于 Jacquet 与 Shalika 对阿基米德
局部积分的扩充一并推广。此事看来没有障碍，但我们没有做。
\begin{annote}
这里点出了本章理论的一处真实缺口：$m<n-1$ 时阿基米德处的有界性没有落实。
读者若只记住定理 4.1 的三条结论，容易以为理论对一切 $m<n$ 都已完整；
实际上 $m<n-1$ 的有界性要绕道 Gelbart--Shahidi 的常数项方法，
或在 $V_\pi\otimes V_{\pi'}$ 上重做一遍积分表示。
原文措辞很谨慎，只说「看来没有障碍」而不说已证，译文保留这一分寸。
\par\vspace{0.15em}{\raggedleft 【来源：原书 p.~221；Gelbart--Shahidi, \textit{Boundedness of automorphic $L$-functions in vertical strips}, J.\ Amer.\ Math.\ Soc.\ 14 (2001), 79--107（A）】\par}
\end{annote}

应当指出，若改用常数项与 Eisenstein 级数 Fourier 系数的方法处理这些
$L$ 函数，则 Gelbart 与 Shahidi 已证明一大类自守 $L$ 函数（包括本章所论者）
在竖直带中有界 [12]。

\subsection{\texorpdfstring{$L$}{L} 函数的极点}

下面确定整体 $L$ 函数可能在何处有极点。$L$ 函数的极点与整体积分的极点相关。
回忆第 2 节：当 $m<n$ 时整体积分 $I(s;\varphi,\varphi')$ 是整函数；
而当 $m=n$ 时，$I(s;\varphi,\varphi',\Phi)$ 至多有单极点，
且仅当 $\pi\cong\widetilde{\pi}'\otimes|\det|^{i\sigma}$（$\sigma$ 实）时，
在 $s=-i\sigma$ 与 $s=1-i\sigma$ 处出现。
如上所述，对适当的 $\varphi$、$\varphi'$ 与 $\Phi$，整体积分与整体 $L$ 函数由
\[
I(s;\varphi,\varphi')=\left(\prod_{v\in S}e_v(s;W_{\varphi_v},W'_{\varphi'_v})\right)
L(s,\pi\times\pi')
\]
或
\[
I(s;\varphi,\varphi',\Phi)
=\left(\prod_{v\in S}e_v(s;W_{\varphi_v},W'_{\varphi'_v},\Phi_v)\right)
L(s,\pi\times\pi')
\]
相联系。

另一方面，对任一 $s_0\in\mathbb{C}$ 与任一 $v$，都可选取局部
$W_v$、$W'_v$，必要时还有 $\Phi_v$，使局部比值
$e_v(s_0;W_v,W'_v)\neq0$ 或 $e_v(s_0;W_v,W'_v,\Phi_v)\neq0$。
于是当 $\varphi$、$\varphi'$ 与 $\Phi$ 在 $v\in S$ 各处变动时，
用这些因子作除法不会给 $L(s,\pi\times\pi')$ 引入多余的极点。
这与 $L$ 因子按局部积分之极点所作的局部刻画相符：
整体地看，$L(s,\pi\times\pi')$ 的极点正是整体积分族
$\{I(s;\varphi,\varphi')\}$ 或 $\{I(s;\varphi,\varphi',\Phi)\}$ 的极点。
故由定理 2.1 与定理 2.2 得下述结论。

\textbf{定理 4.2.}\quad 若 $m<n$，则 $L(s,\pi\times\pi')$ 是整函数。
若 $m=n$，则 $L(s,\pi\times\pi')$ 至多有单极点，
且仅当 $\pi\cong\widetilde{\pi}'\otimes|\det|^{i\sigma}$（$\sigma$ 实）时出现，
位置为 $s=-i\sigma$ 与 $s=1-i\sigma$。

由此可得两条有用的观察。
\begin{enumerate}
\item $L(s,\pi\times\widetilde{\pi})$ 在 $s=0$ 与 $s=1$ 处有单极点。
\item 对 $\GL_n(\A)$ 的尖点表示 $\pi$ 与 $\pi'$，$L(s,\pi\times\widetilde{\pi}')$
在 $s=1$ 处有极点，当且仅当 $\pi\cong\pi'$。
\end{enumerate}
作为推论，我们得到强重数一定理的解析证明 [31, 25]。

\textbf{强重数一.}\quad 设 $(\pi,V_\pi)$ 与 $(\pi',V_{\pi'})$ 是 $\GL_n(\A)$ 的
两个尖点表示。若存在位置的有限集 $S$，使对一切 $v\notin S$ 有
$\pi_v\cong\pi'_v$，则 $\pi=\pi'$。
\begin{annote}
把 $L(s,\pi\times\widetilde{\pi}')$ 在 $s=1$ 的极点当作 $\pi\cong\pi'$ 的判别，
是 Rankin--Selberg 理论最漂亮的应用之一。
逻辑上它依赖两件事：一是 $L$ 因子在 $s=1$ 的留数可以写成
$\pi$ 与 $\pi'$ 的内积（这正是第 2.3 节中当作「不合适的 $L^2$ 内积」
而暂时搁置的那个积分），二是定理 4.2 保证除这一处外没有别的极点。
于是「几乎所有局部分量相同」立刻升级为「整体同构」，
这正是强重数一，而重数一只是它在 $S=\varnothing$ 时的特例。
\par\vspace{0.15em}{\raggedleft 【来源：Jacquet--Shalika, Amer.\ J.\ Math.\ 103 (1981), 777--815；Piatetski-Shapiro, \textit{Multiplicity one theorems}, Proc.\ Sympos.\ Pure Math.\ 33 (1979)（A）】\par}
\end{annote}
\section{逆定理}

让我们先回到 Hecke。回忆对一个模形式
\[
f(\tau)=\sum_{n=1}^{\infty}a_ne^{2\pi in\tau},
\]
例如关于 $\mathrm{SL}_2(\mathbb{Z})$ 的，Hecke 配上 $L$ 函数 $L(s,f)$，
二者由 Mellin 变换相联系：
\[
\Lambda(s,f)=(2\pi)^{-s}\Gamma(s)L(s,f)=\int_0^\infty f(iy)y^s\,d^\times y,
\]
而 $L(s,f)$ 的函数方程由 $f(\tau)$ 在 $\tau\mapsto-1/\tau$ 下的模变换律导出。
在其奠基性论文 [17] 中，Hecke 把这一过程倒了过来。
他取一个 Dirichlet 级数
\[
D(s)=\sum_{n=1}^{\infty}\frac{a_n}{n^s},
\]
假定它在某半平面收敛，可整延拓为有限阶函数，并满足与权 $k$ 的模形式的 $L$ 函数
相同的函数方程；于是实际上可由 $D(s)$ 经 Mellin 逆变换重构出模形式
\[
f(iy)=\sum_na_ne^{-2\pi ny}
=\frac{1}{2\pi i}\int_{2-i\infty}^{2+i\infty}(2\pi)^{-s}\Gamma(s)D(s)y^s\,ds,
\]
并由 $D(s)$ 在 $s\mapsto k-s$ 下的函数方程，得到 $f(\tau)$ 在
$\tau\mapsto-1/\tau$ 下的模变换律。这就是 Hecke 逆定理。

本节将给出 Hecke 定理在 $\GL_n$ 的 $L$ 函数框架中的一些类比。
令人惊讶的是，技巧与 Hecke 完全相同，即把积分表示反过来解。
此事最先在 $\GL_2$ 的自守表示框架中由 Jacquet 与 Langlands 完成 [21]，
随后由 Jacquet、Piatetski-Shapiro 与 Shalika 对 $\GL_3$ 推广并显著加强 [22]。
本节主要取自我们的综述 [7]，进一步的细节见 [5, 6]。

设 $k$ 为整体域，$\A$ 为其阿代尔环，$\psi$ 为 $\A$ 的一个在 $k$ 上平凡的
非平凡连续加法特征。以下取整数 $n\ge3$。

为陈述这些逆定理，我们从一个 $\GL_n(\A)$ 的不可约容许表示 $\Pi$ 出发。
按本注记的惯例，假定 $\Pi$ 酉且一般，但这一假定并不必要。
它有分解 $\Pi=\otimes'\Pi_v$，其中 $\Pi_v$ 是 $\GL_n(k_v)$ 的不可约容许一般表示。
由第 3 节的局部理论，每个 $\Pi_v$ 配有局部 $L$ 函数 $L(s,\Pi_v)$ 与
局部 $\varepsilon$ 因子 $\varepsilon(s,\Pi_v,\psi_v)$。故形式上可作
\[
L(s,\Pi)=\prod_vL(s,\Pi_v),
\qquad
\varepsilon(s,\Pi,\psi)=\prod_v\varepsilon(s,\Pi_v,\psi_v).
\]
我们始终假定 $\Pi$ 满足以下两件事：
\begin{enumerate}
\item $L(s,\Pi)$ 在某半平面 $\operatorname{Re}(s)\gg0$ 收敛；
\item $\Pi$ 的中心特征 $\omega_\Pi$ 是自守的，即在 $k^\times$ 下不变。
\end{enumerate}
在这些假定下，$\varepsilon(s,\Pi,\psi)=\varepsilon(s,\Pi)$ 与 $\psi$ 的选取无关 [5]。

我们的逆定理将涉及用 $\GL_m(\A)$ 的尖点自守表示作扭，$m$ 取某些值。
为方便计，记 $\mathcal{A}(m)$ 为 $\GL_m(\A)$ 的自守表示之集，
$\mathcal{A}_0(m)$ 为 $\GL_m(\A)$ 的尖点表示之集，并置
$\mathcal{T}(m)=\bigcup_{d=1}^m\mathcal{A}_0(d)$。
若固定一个由有限位置组成的有限集 $S$，则以 $\mathcal{T}(S;m)$ 记
$\mathcal{T}(m)$ 中在所有 $v\notin S$ 处无分歧的表示所成子集。

设 $\pi'=\otimes'\pi'_v$ 是 $\GL_m(\A)$ 的尖点表示，$m<n$。于是又可形式地定义
\[
L(s,\Pi\times\pi')=\prod_vL(s,\Pi_v\times\pi'_v),
\qquad
\varepsilon(s,\Pi\times\pi')=\prod_v\varepsilon(s,\Pi_v\times\pi'_v,\psi_v),
\]
因为局部因子无论 $\Pi$ 是否自守都有意义。
由上两条假定及 $\pi'$ 的尖点性，$L(s,\Pi\times\pi')$ 与
$L(s,\widetilde{\Pi}\times\widetilde{\pi}')$ 都在 $\operatorname{Re}(s)\gg0$
绝对收敛，且 $\varepsilon(s,\Pi\times\pi')$ 与 $\psi$ 的选取无关。

若 $L(s,\Pi\times\pi')$ 具有 $\Pi$ 为尖点时它应有的解析性质，
则称它\emph{好的}（nice），即：
\begin{enumerate}
\item $L(s,\Pi\times\pi')$ 与 $L(s,\widetilde{\Pi}\times\widetilde{\pi}')$
可解析延拓为 $s$ 的整函数；
\item 这些整延拓在有限宽的竖直带中有界；
\item 它们满足标准函数方程
$L(s,\Pi\times\pi')=\varepsilon(s,\Pi\times\pi')L(1-s,\widetilde{\Pi}\times\widetilde{\pi}')$。
\end{enumerate}

$\GL_n$ 的基本逆定理如下 [5, 4]。

\textbf{定理 5.1.}\quad 设 $\Pi$ 是如上 $\GL_n(\A)$ 的不可约容许表示，
$S$ 是有限位置的有限集。若对一切 $\pi'\in\mathcal{T}(S;n-1)$，
$L(s,\Pi\times\pi')$ 都是好的，则 $\Pi$ 是拟自守的，
意即存在自守表示 $\Pi'$，使对一切 $v\notin S$ 有 $\Pi_v\cong\Pi'_v$。
若 $S$ 为空，则实际上 $\Pi$ 是 $\GL_n(\A)$ 的尖点自守表示。

此结果对 $n=2$ 当然也成立。

为应用 [3]，希望扭得越少越好。限制扭的方式本质上只有两种：
一是限制扭表示所在群的秩，二是限制分歧。
\begin{annote}
逆定理的逻辑值得单独说清，因为它与「正定理」的方向相反。
正向理论从自守表示出发造出 $L$ 函数，并证明它满足函数方程与整性；
逆定理则把这些性质当作公理，反推出 $\Pi$ 必为自守。
Hecke 用一个 Mellin 逆变换完成反推，$\GL_n$ 上没有现成的逆变换，
替代物是「反解积分表示」：先假定 $L(s,\Pi\times\pi')$ 好，
于是对每个试验表示 $\pi'$ 都能把积分表示反解出 $\Pi$ 的 Whittaker 函数，
再验证这些 Whittaker 函数确来自一个自守表示。
扭集取 $\mathcal{T}(S;n-1)$ 而非全部尖点表示，是因为证明中只需这些扭。
\par\vspace{0.15em}{\raggedleft 【来源：Cogdell--Piatetski-Shapiro, \textit{Converse theorems for $\mathrm{GL}_n$}, Publ.\ Math.\ IH\'ES 79 (1994), 157--214（A）】\par}
\end{annote}

限制扭的秩，可得下述结果 [6]。

\textbf{定理 5.2.}\quad 设 $\Pi$ 是如上 $\GL_n(\A)$ 的不可约容许表示，
$S$ 是有限位置的有限集。若对一切 $\pi'\in\mathcal{T}(S;n-2)$，
$L(s,\Pi\times\pi')$ 都是好的，则 $\Pi$ 是拟自守的，
意即存在自守表示 $\Pi'$，使对一切 $v\notin S$ 有 $\Pi_v\cong\Pi'_v$。
若 $S$ 为空，则实际上 $\Pi$ 是 $\GL_n(\A)$ 的尖点自守表示。

此结果比定理 5.1 强，但证明更精细。

为应用这些定理，须把集合 $S$ 的选取与用充分分歧的特征 $\eta$ 作扭配合起来。
集合 $S$ 通常由 $\Pi_v$ 分歧的那些有限位置组成。
$\eta$ 的用法则与前面提到的 $\gamma$ 因子的局部稳定性相配合。
此时下述观察是应用上述任一定理的关键 [7]。

\textbf{观察.}\quad 设 $\Pi$ 如定理 5.1 或 5.2。设 $\eta$ 是 $k^\times\backslash\A^\times$
的一个固定（充分分歧）特征。若对一切 $\pi'\in\mathcal{T}\otimes\eta$，
$L(s,\Pi\times\pi')$ 都是好的，其中 $\mathcal{T}$ 是定理 5.1 或 5.2 的任一扭集，
则 $\Pi$ 如那些定理中所述是拟自守的。

只需注意（例如看局部或整体积分即知）：若 $\pi'\in\mathcal{T}$，则
$L(s,\Pi\times(\pi'\otimes\eta))=L(s,(\Pi\otimes\eta)\times\pi')$。
于是对 $\Pi$ 用扭集 $\mathcal{T}\otimes\eta$ 应用逆定理，
等价于对 $\Pi\otimes\eta$ 用扭集 $\mathcal{T}$ 应用逆定理。
故由定理 5.1 或 5.2（视何者适用），$\Pi\otimes\eta$ 拟自守，从而 $\Pi$ 亦然。

限制扭的第二种方式，是除有限多个位置外限制分歧 [5]。
现在固定一个非空的有限位置集 $S$；在数域情形还要求 $S$ 含全体阿基米德位置
$S_\infty$。以 $\mathcal{T}_S(m)$ 记 $\mathcal{T}(m)$ 中满足「对一切 $v\notin S$
无分歧」的表示 $\pi'$ 所成子集。
注意这是对这些表示的分歧加了极重的限制。

\textbf{定理 5.3.}\quad 设 $\Pi$ 是如上 $\GL_n(\A)$ 的不可约容许表示，
$S$ 是非空的有限位置集，含 $S_\infty$，且使 $S$-整数环 $\mathfrak{o}_S$ 的类数为一。
若对一切 $\pi'\in\mathcal{T}_S(n-1)$，$L(s,\Pi\times\pi')$ 都是好的，
则 $\Pi$ 是拟自守的，意即存在自守表示 $\Pi'$，使对一切 $v\in S$ 有
$\Pi_v\cong\Pi'_v$，且对一切 $v\notin S$（$\Pi_v$ 与 $\Pi'_v$ 均无分歧者）亦然。

此处有几点须注意。第一，这里出现了类数限制。
不过若 $k=\mathbb{Q}$，则可取 $S=S_\infty$，
于是得到只用「level 1」扭的逆定理。
从实用角度考虑，若以 $S_\Pi$ 记 $\Pi_v$ 分歧的有限位置之集，
则在应用中通常取 $S$ 与 $S_\Pi$ 不相交。
此时凡 $v\notin S$ 处我们又丢失了全部信息，除非 $\Pi_v$ 本就在该处无分歧。
\begin{annote}
类数条件看似技术性，实则是这一版逆定理的命门。
当 $S=S_\infty$ 时，$\mathcal{T}_S(n-1)$ 里只有 level 1 的尖点表示，
可用的扭少到极点；类数为一正是让「level 1」的 Hecke 特征足够多、
从而能用强逼近把局部条件拼成整体结论的临界条件。
$k=\mathbb{Q}$ 时 $\mathfrak{o}_S=\mathbb{Z}$ 类数为一，故定理涵盖经典情形。
若类数大于一，须把扭扩大到某个射线类群的特征，定理形式随之改变。
\par\vspace{0.15em}{\raggedleft 【来源：Cogdell--Piatetski-Shapiro, \textit{Converse theorems for $\mathrm{GL}_n$, II}, J.\ reine angew.\ Math.\ 507 (1999), 165--188（A）】\par}
\end{annote}

$S$ 为空时定理 5.1 的证明，大体沿 Hecke、Weil 与 Jacquet--Langlands 的路子。
它基于 $L$ 函数的积分表示、Fourier 展开、Mellin 逆变换，
最后用到 Langlands 谱理论的弱形式。
对定理 5.1、5.2、5.3，由于限制了扭，须施加若干局部条件以补偿扭的不足。
定理 5.1 与 5.2 只有有限多个局部条件，定理 5.3 则有无限多个。
于是须借助同余子群的生成以及弱逼近或强逼近的结果来绕过这些条件。

\begin{annote}
底本笔误一览，译文均照录原文，读者按此校正。
(a) 原书 p.~201 第 1.2 节首段「for all $n\in N_k(k_v)$」，其中 $N_k$ 应为 $N_n$。
(b) 原书 p.~211 称 $q=|\mathfrak{o}/\mathfrak{p}|$ 为「剩余域次数」，实为剩余域基数。
(c) 原书 p.~213 局部 $L$ 函数一节中写「the representation of $\GL_n(\A)$」，
此处的群应为 $\GL_n(k)$，与全节的局部设定一致。
(d) 原书 p.~218 的 $\Psi(s;W,\Phi)$ 定义末尾写作 $dh$，积分变量应为 $dg$。
(e) 原书 p.~210「Hence a residues can be nonzero」衍一冠词，应为「the residues」。
(f) 原书 p.~214「except for a finite number if exceptional values」，
其中 if 应为 of。
(g) 原书 p.~212 把 $\mathcal{S}(\A^n)$ 上的 Fourier 变换写成对 $\A^\times$ 积分，
积分区域应为 $\A^n$。
以上七处均经页图逐字核对，不涉数学内容，仅涉记号与拼写。
\par\vspace{0.15em}{\raggedleft 【来源：原书 pp.~201, 210, 211, 212, 213, 214, 218 影印页；\mbox{底本 PDF 204, 213, 214, 215, 216, 217, 221}（A）】\par}
\end{annote}

\par\medskip
\noindent\textit{（参考文献 49 条见原书 pp.~225--228，即底本 PDF 228--231；限于篇幅不逐条迻录。）}
